% concatenated from sn-coleman.tex
%Version 3.1 December 2024
% See section 11 of the User Manual for version history
%
%%%%%%%%%%%%%%%%%%%%%%%%%%%%%%%%%%%%%%%%%%%%%%%%%%%%%%%%%%%%%%%%%%%%%%
%%                                                                 %%
%% Please do not use \input{...} to include other tex files.       %%
%% Submit your LaTeX manuscript as one .tex document.              %%
%%                                                                 %%
%% All additional figures and files should be attached             %%
%% separately and not embedded in the \TeX\ document itself.       %%
%%                                                                 %%
%%%%%%%%%%%%%%%%%%%%%%%%%%%%%%%%%%%%%%%%%%%%%%%%%%%%%%%%%%%%%%%%%%%%%

%%\documentclass[referee,sn-basic]{sn-jnl}% referee option is meant for double line spacing

%%=======================================================%%
%% to print line numbers in the margin use lineno option %%
%%=======================================================%%

%%\documentclass[lineno,pdflatex,sn-basic]{sn-jnl}% Basic Springer Nature Reference Style/Chemistry Reference Style

%%=========================================================================================%%
%% the documentclass is set to pdflatex as default. You can delete it if not appropriate.  %%
%%=========================================================================================%%

%%\documentclass[sn-basic]{sn-jnl}% Basic Springer Nature Reference Style/Chemistry Reference Style

%%Note: the following reference styles support Namedate and Numbered referencing. By default the style follows the most common style. To switch between the options you can add or remove “Numbered” in the optional parenthesis. 
%%The option is available for: sn-basic.bst, sn-chicago.bst%  
 
%%\documentclass[pdflatex,sn-nature]{sn-jnl}% Style for submissions to Nature Portfolio journals
%%\documentclass[pdflatex,sn-basic]{sn-jnl}% Basic Springer Nature Reference Style/Chemistry Reference Style
\documentclass[pdflatex,sn-mathphys-num]{sn-jnl}% Math and Physical Sciences Numbered Reference Style
%%\documentclass[pdflatex,sn-mathphys-ay]{sn-jnl}% Math and Physical Sciences Author Year Reference Style
%%\documentclass[pdflatex,sn-aps]{sn-jnl}% American Physical Society (APS) Reference Style
%%\documentclass[pdflatex,sn-vancouver-num]{sn-jnl}% Vancouver Numbered Reference Style
%%\documentclass[pdflatex,sn-vancouver-ay]{sn-jnl}% Vancouver Author Year Reference Style
%%\documentclass[pdflatex,sn-apa]{sn-jnl}% APA Reference Style
%%\documentclass[pdflatex,sn-chicago]{sn-jnl}% Chicago-based Humanities Reference Style

%%%% Standard Packages
%%<additional latex packages if required can be included here>
\usepackage[utf8]{inputenc}
\usepackage[T1]{fontenc}
\usepackage{graphicx}%
\usepackage{multirow}%
\usepackage{amsmath,amssymb,amsfonts}%
\usepackage{amsthm}%
\usepackage{mathrsfs}%
\usepackage[title]{appendix}%
\usepackage{xcolor}%
\usepackage{textcomp}%
\usepackage{manyfoot}%
\usepackage{booktabs}%
\usepackage{algorithm}%
\usepackage{algorithmicx}%
\usepackage{algpseudocode}%
\usepackage{listings}%
\usepackage{cleveref}
\usepackage{tablefootnote}
\usepackage{enumitem}

% Math shortcuts
\DeclareMathOperator*{\argmin}{arg\,min}
\DeclareMathOperator*{\argmax}{arg\,max}
\DeclareMathOperator{\Tr}{tr}
\newcommand{\R}{\mathbb{R}}
\newcommand{\E}{\mathbb{E}}
\newcommand{\norm}[2]{\|#1\|_{#2}}
\newcommand{\Anorm}[1]{\|#1\|_{A}}
\newcommand{\lmin}{\lambda_{\min}}
\newcommand{\lmax}{\lambda_{\max}}

%% as per the requirement new theorem styles can be included as shown below
\theoremstyle{thmstyleone}%
\newtheorem{theorem}{Theorem}%  meant for continuous numbers
%%\newtheorem{theorem}{Theorem}[section]% meant for sectionwise numbers
%% optional argument [theorem] produces theorem numbering sequence instead of independent numbers for Proposition
\newtheorem{proposition}[theorem]{Proposition}% 
%%\newtheorem{proposition}{Proposition}% to get separate numbers for theorem and proposition etc.

% Theorem environments

\newtheorem{corollary}[theorem]{Corollary}
\newtheorem{lemma}[theorem]{Lemma}


\theoremstyle{thmstyletwo}%
\newtheorem{example}{Example}%
\newtheorem{remark}{Remark}%

\theoremstyle{thmstylethree}%
\newtheorem{definition}{Definition}%

\raggedbottom
%%\unnumbered% uncomment this for unnumbered level heads

\begin{document}
\title[Residual-Weighted Randomized Jacobi]{Residual-Weighted Randomized Jacobi: Sharpened Bounds via Residual Concentration and Asynchronous Extension}




%%=============================================================%%
%% GivenName	-> \fnm{Joergen W.}
%% Particle	-> \spfx{van der} -> surname prefix
%% FamilyName	-> \sur{Ploeg}
%% Suffix	-> \sfx{IV}
%% \author*[1,2]{\fnm{Joergen W.} \spfx{van der} \sur{Ploeg} 
%%  \sfx{IV}}\email{iauthor@gmail.com}
%%=============================================================%%

\author*[1]{\fnm{Evan} \sur{Coleman}}\email{ecolema4@umw.edu}


\affil*[1]{\orgdiv{Department of Computer Science}, \orgname{University of Mary Washington}, \city{Fredericksburg},  \state{Virginia}, \country{USA}}


%%==================================%%
%% Sample for unstructured abstract %%
%%==================================%%

\abstract{We study randomized stationary methods for symmetric positive definite
linear systems in which component $j$ is selected with probability
proportional to $|r_j|^\ell$. This power-weighted family interpolates
continuously between uniform randomized Jacobi as $\ell \to 0$ and
Gauss--Southwell greedy relaxation as $\ell \to \infty$.
For the central case $\ell = 2$, we sharpen the standard one-step
convergence analysis using the inverse participation ratio (IPR)
$\nu^2(r) = n\|r\|_4^4/\|r\|_2^4$, which equals $1$ when the residual
is uniform and grows toward $n$ as it concentrates. The resulting
bound amplifies the expected per-step progress by exactly $\nu^2$
over the uniform-sampling baseline. The IPR
can be computed online at $O(n)$ cost and doubles as a per-iteration
diagnostic.

We extend the analysis to asynchronous power-weighted Jacobi via the
Avron--Druinsky--Gupta framework, obtaining an epoch-based convergence
theorem in which the IPR controls both the progress coefficient and
the allowed-delay window.
Numerical experiments on shared-memory hardware support the sharpened
bound and show the IPR trajectory is essentially concurrency-insensitive.
Unexpectedly, consistent-reads execution, the easier case for the ADG
analysis, destabilizes power-weighted sampling at high concurrency
while inconsistent reads remain stable; the same IPR that amplifies
progress 
amplifies a thread-collision rate that inconsistent reads appear to absorb.
We propose a feedback-damping mechanism and verify
two predictions about its dependence on problem size.}





\keywords{Asynchronous iterative methods, Randomized Jacobi, Inverse participation ratio, Residual-weighted sampling}

%%\pacs[JEL Classification]{D8, H51}

%%\pacs[MSC Classification]{35A01, 65L10, 65L12, 65L20, 65L70}

\maketitle

\section{Introduction}
\label{sec:intro}

Randomized iterative methods for solving linear systems have become a standard
tool in numerical linear algebra, optimization, and scientific computing.  In
row-action and coordinate-relaxation methods, each iteration updates only one
row, coordinate, component, or block of the current approximation.  Randomized
Kaczmarz methods \cite{strohmer2009randomized}, randomized coordinate descent,
and randomized Gauss--Seidel-type schemes \cite{leventhal2010randomized}
illustrate the appeal of this approach: each step is inexpensive, the update
formula is simple, and the method can often be implemented with relatively low
synchronization overhead.  These properties make randomized stationary methods
especially natural candidates for asynchronous parallel computation
\cite{avron2015revisiting}.

A central design choice in such methods is the component-selection rule.  The
simplest choice is uniform random sampling, which is easy to implement and
analyze.  More refined fixed distributions, such as row-norm or
diagonal-weighted sampling, incorporate matrix geometry but remain independent
of the current residual.  At the opposite extreme are greedy rules such as
Southwell relaxation \cite{southwell1946relaxation}, which
selects the component with the largest residual or largest predicted progress.
Greedy selection can be substantially more efficient per iteration, but it
requires global residual comparisons, sorting, or priority maintenance.  These
operations can dominate the cost of a relaxation step and are particularly
awkward in asynchronous or distributed settings.

This paper studies a middle ground between fixed randomized sampling and fully
greedy relaxation.  Given the current residual $r^{(i)}=b-Ax^{(i)}$, we select
component $j$ with probability
\[
    P(k=j) =
    \frac{|r_j^{(i)}|^\ell}{\sum_m |r_m^{(i)}|^\ell},
    \qquad \ell>0.
\]
This \emph{power-weighted} family is residual-adaptive but does not require
sorting, thresholding, or constructing an eligible set.  It also has transparent
limiting behavior: as $\ell \to 0$, the method approaches uniform randomized
Jacobi, while as $\ell \to \infty$, it approaches Gauss--Southwell selection.  The
parameter $\ell$ therefore controls how aggressively the method concentrates
sampling effort on large-residual components.

The motivation for this family is simple.  Uniform sampling treats all
components as equally important, even when most of the residual mass is
concentrated in only a few entries.  In such a regime, many uniform updates are
spent on components that can contribute little immediate progress.  Residual-weighted sampling instead increases the probability of selecting components
that currently dominate the residual.  The key question is whether this
intuition can be made quantitative.  In particular, one would like a convergence
bound that reduces to the uniform bound when the residual is flat, but becomes
strictly sharper when the residual is non-uniform.

For synchronous randomized Jacobi, we answer this question using the
inverse participation ratio (IPR) of the residual,
\[
\nu^2(r) = \frac{n \cdot \norm{r}{4}^4}{\norm{r}{2}^4},
\]
a scalar that equals one when all residual components have equal magnitude
and grows toward $n$ as the residual concentrates on a single component.
For power-weighted sampling with $\ell = 2$, the expected one-step progress
naturally contains a residual-norm ratio that rewrites as a multiple of the
IPR, and the resulting sharpened bound amplifies the expected per-step
progress by exactly that scalar. The same quantity that explains \emph{when}
adaptive sampling helps therefore quantifies \emph{how much} it helps.

The asynchronous setting is more delicate. In uniform asynchronous Jacobi,
the component choice is independent of the stale iterate, and the expected
interference between simultaneously updated components is averaged away by
the uniform sampling distribution. Residual-weighted sampling breaks both
features: the sampling probabilities depend on the stale residual seen by
the worker, coupling the choice to the delay, and concentrated residuals
produce concentrated sampling distributions whose interference is no longer
smoothed by the averaging the uniform analysis relies on. The result is a
genuine tradeoff: residual concentration increases expected progress but
can also increase the worst-case interference bound. We extend the
Avron--Druinsky--Gupta (ADG) framework for uniform asynchronous
Jacobi~\cite{avron2015revisiting} to this residual-dependent setting under
a uniform bound on the IPR trajectory. The resulting theorem is structurally
analogous to ADG's epoch-reduction theorem, with the IPR controlling both
the progress coefficient and the allowed-delay window, and should be read
as a first asynchronous convergence result for power-weighted sampling
rather than as a strict generalization of the uniform theory.

The IPR sits well above its uniform floor of 1 on every test problem
in our suite, with steady-state values of $\nu^2 \in [4, 6]$ across the
three problem classes and substantially higher early-iteration values on
instances with localized forcing. Hardware experiments on a 96-core EPYC
node confirm that power-weighted sampling translates the IPR advantage
into measurably faster convergence under emergent asynchrony, and the IPR
trajectory is essentially insensitive to thread count: concurrency-induced
variation stays under 3\% across two decades of thread count, supporting
the global trajectory bound used in the asynchronous theorem below.

The contributions of this paper are as follows.

\begin{enumerate}
    \item We organize randomized, adaptive, and greedy relaxation methods into a
    component-selection adaptivity spectrum, positioning residual-weighted
    sampling between fixed randomized methods and Southwell-type greedy
    relaxation.

    \item We prove that power-weighted Jacobi with any $\ell>0$ matches the
    standard uniform randomized Jacobi convergence bound in the worst case.

    \item For $\ell=2$, we derive an IPR-sharpened one-step convergence bound showing
    that the expected progress is amplified by the residual's inverse
    participation ratio.  This gives a theoretical separation between uniform
    and residual-weighted sampling whenever the residual is non-uniform.

    \item We extend the analysis to asynchronous power-weighted Jacobi.  Under a
    uniform IPR-trajectory assumption, we obtain an ADG-style epoch reduction
    theorem that exposes the progress--interference tradeoff created by
    residual-dependent sampling.

    \item We investigate the IPR and convergence behavior numerically on
    several example problems, with the goal of assessing whether the residual
    concentration required by the sharpened bounds occurs in practical
    stationary iterations and whether it persists under asynchronous execution.

    \item We observe empirically that the theorem above assumes consistent
    reads (the standard easier case in the ADG framework), but our hardware
    experiments reveal that consistent read execution is in fact the unsafe case for
    power-weighted sampling at high concurrency: consistent reads destabilize
    the iteration on all three test problems, while inconsistent reads remain
    stable. %We conjecture a feedback-damping mechanism in which the same IPR
    %that drives progress amplification also drives a collision rate that
    %only inconsistent reads can absorb, and we identify a sign-aware
    %cross-term bound as the natural next theoretical step.
\end{enumerate}

The rest of the paper is organized as follows.  \Cref{sec:related-work}
reviews asynchronous iterative methods and related randomized and adaptive
selection rules.  \Cref{sec:background} introduces the randomized Jacobi
framework, the one-step progress identity, baseline convergence bounds, and the
component-selection adaptivity spectrum.  \Cref{sect:baseline-residual-weighted}
develops the synchronous residual-weighted analysis, including the worst-case
power-weighted bound and the IPR-sharpened result.  \Cref{sec:async} extends the
analysis to asynchronous power-weighted Jacobi and identifies the main sources
of looseness in the current theorem.  \Cref{sec:experiments} presents numerical
experiments, and \Cref{sec:conclusion} concludes.


\section{Related Work}
\label{sec:related-work}

Asynchronous iterative methods have a long history in numerical linear algebra,
beginning with the chaotic relaxation framework of Chazan and Miranker
\cite{chazan1969chaotic} and the multiprocessor analysis of Baudet
\cite{baudet1978asynchronous}.  These early works established that classical
stationary iterations can converge even when processors update components using
stale information, provided the delays and communication patterns satisfy
appropriate regularity assumptions.  The broader theory of asynchronous parallel
computation was later developed systematically by Bertsekas and Tsitsiklis
\cite{bertsekas1989parallel}, and the mathematical foundations and convergence
conditions for asynchronous iterations were surveyed by Frommer and Szyld
\cite{frommer2000asynchronous}.  In this classical literature, the primary goal
is to characterize when asynchrony preserves convergence; sharp convergence
rates and adaptive component-selection rules are typically secondary concerns.

A probabilistic viewpoint on asynchronous iteration was developed by Strikwerda
\cite{strikwerda2002probabilistic}, who extended deterministic asynchronous
models by allowing both component choices and delays to be described
probabilistically.  This perspective is particularly relevant for modern
parallel and distributed implementations, where update order, communication
delays, and processor availability may be irregular or random.  More recently,
Avron, Druinsky, and Gupta revisited asynchronous linear solvers for symmetric
positive definite systems and provided convergence-rate bounds for randomized
asynchronous Jacobi under bounded-delay and independence assumptions
\cite{avron2015revisiting}.  Their work is the closest theoretical point of
comparison for the asynchronous analysis in this paper: like their model, we
study randomized component updates applied to SPD systems, but we replace fixed
uniform sampling with residual-dependent sampling and therefore must account for
the coupling between stale residuals, sampling probabilities, and delayed
updates.

Asynchronous stationary iterations are also of practical interest because they
can reduce synchronization overhead in settings where exact high-accuracy solves
are not required at every stage.  Chow and Patel developed fine-grained parallel
algorithms for incomplete LU factorization, motivated by the use of incomplete
factorizations as scalable preconditioners for linear solvers
\cite{chow2015fine}.  Related work has studied asynchronous Jacobi iterations
both as standalone solvers and as building blocks for larger methods.  For
example, Wolfson-Pou and Chow developed convergence models for asynchronous
Jacobi and observed behavior that can differ substantially from simple delay
models \cite{wolfsonpou2018convergence}.  They also investigated asynchronous
multigrid methods, where asynchronous relaxation can act as a smoother and help
reduce synchronization costs within multilevel solvers
\cite{wolfsonpou2019multigrid}.  Chow, Frommer, and Szyld analyzed
asynchronous Richardson iterations, further clarifying the relationship between
classical stationary methods, relaxation parameters, and asynchronous execution
\cite{chow2021asynchronous}.  More recent work has extended these ideas to
asynchronous semi-iterative and Chebyshev-type methods, including multigrid
preconditioning \cite{wolfsonpou2025asynchronous}.

Application-oriented studies have shown asynchronous linear iterations to be
useful in large-scale PDE contexts~\cite{glusa2020scalable,kashi2021asynchronous,kashi2025effectiveness},
particularly as low-synchronization components within preconditioners,
smoothers, and domain-decomposition solvers rather than as replacements for
Krylov methods in high-accuracy regimes.
The same synchronization--staleness tradeoff appears in neighboring areas:
asynchronous coordinate-update methods such as HOGWILD!~\cite{recht2011hogwild},
asynchronous stochastic coordinate descent~\cite{liu2015asynchronous}, and the
ARock fixed-point framework~\cite{peng2016arock} all exploit it in optimization
and machine learning, providing useful context even though the present paper
focuses on stationary solvers for linear systems.
Finally, asynchronous methods are closely connected to questions of resilience
and irregular execution.  In distributed or cloud-based environments, workers
may be delayed, temporarily unavailable, or subject to heterogeneous performance.
Recent straggler-tolerant stationary methods exploit this observation by
allowing iterations to proceed using partial or delayed information
\cite{kalantzis2025straggler}.  Fault-tolerant variants of asynchronous linear
solvers have also been studied in the context of data corruption and recovery
\cite{coleman2021fault,coleman2026fault,vogl2024modifying}.  These directions are not the focus
of the present paper, but they reinforce the broader motivation for analyzing
stationary methods under stale and irregular update patterns.

The present work fits into this literature by studying how adaptive,
residual-dependent component selection interacts with asynchronous stationary
iteration.  Classical asynchronous theory explains when stale updates can be
tolerated; randomized asynchronous Jacobi provides a rate baseline under uniform
sampling; and practical studies motivate the use of asynchronous stationary
iterations as scalable smoothers, preconditioners, and low-accuracy solvers.  Our
contribution is complementary: we ask whether the component-selection rule itself
can be biased toward high-residual entries while retaining useful convergence
guarantees, and we quantify the resulting advantage through residual
non-uniformity.


\section{Background}
\label{sec:background}

\subsection{Setup and Notation}

Let $A \in \R^{n \times n}$ be symmetric positive definite (SPD) and let
$b \in \R^n$.  We consider the linear system
\[
    Ax = b,
\]
whose unique solution is denoted by $x^* = A^{-1}b$.  Since $A$ is SPD, it has
real positive eigenvalues
\[
    0 < \lambda_1 \leq \lambda_2 \leq \cdots \leq \lambda_n,
\]
and we write
\[
    \lmin(A) = \lambda_1, \qquad
    \lmax(A) = \lambda_n, \qquad
    \kappa = \frac{\lmax(A)}{\lmin(A)}.
\]
When the matrix is clear from context, we write simply $\lmin$ and $\lmax$.

At iteration $i$, the current approximate solution is $x^{(i)}$.  We denote the
error and residual by
\[
    e^{(i)} = x^{(i)} - x^*, \qquad
    r^{(i)} = b - A x^{(i)}.
\]
Since $Ax^*=b$, the residual and error are related by
\[
    r^{(i)} = -A e^{(i)}.
\]
We measure error in the energy norm induced by $A$,
\[
    \Anorm{z}^2 = z^T A z.
\]
We also use the standard basis vectors $e_1,\ldots,e_n$, where $e_j$ denotes
the vector with a $1$ in component $j$ and zeros elsewhere.  To avoid confusing
the basis vector $e_j$ with the iteration error $e^{(i)}$, the iteration index
will always be written in parentheses.  The diagonal entries of $A$ are denoted
by $A_{jj}$, and we define
\[
    A_{\max} = \max_{1 \leq j \leq n} A_{jj}, \qquad
    A_{\min} = \min_{1 \leq j \leq n} A_{jj}.
\]
For SPD matrices, $A_{jj} = e_j^T A e_j > 0$ for every $j$, so all coordinate
relaxations below are well-defined. We keep the diagonal entries explicit
throughout the synchronous analysis because the baseline sampling
distributions depend on them, even though one could assume $A_{jj} = 1$
without loss of generality after a diagonal change of variables; this
keeps residual-weighted sampling directly comparable to both uniform and
diagonal-weighted sampling.

We will repeatedly use the spectral inequalities
\[
    \lmin \norm{z}{2}^2
    \leq
    \Anorm{z}^2
    \leq
    \lmax \norm{z}{2}^2,
\]
and, since $r=-Ae$,
\begin{equation}
\label{eq:residual-energy-bounds}
    \lmin \Anorm{e}^2
    \leq
    \norm{r}{2}^2
    \leq
    \lmax \Anorm{e}^2.
\end{equation}
The lower bound in~\eqref{eq:residual-energy-bounds} is the main bridge between
expected residual reduction and convergence in the $A$-norm.


\subsection{Randomized Jacobi Relaxation}

A single Jacobi relaxation updates one component of $x$ using the current
residual.  If component $k$ is selected at iteration $i$, the update is
\begin{equation}
\label{eq:update}
    x^{(i+1)}
    =
    x^{(i)} + \frac{r_k^{(i)}}{A_{kk}}\, e_k .
\end{equation}
The effect of one coordinate relaxation on the $A$-norm error is exact.  Since
$r_k^{(i)} = -e_k^T A e^{(i)}$, the new error is
\[
    e^{(i+1)}
    =
    e^{(i)} + \frac{r_k^{(i)}}{A_{kk}} e_k .
\]
Expanding the $A$-norm gives
\begin{align}
    \Anorm{e^{(i+1)}}^2
    &=
    \left(e^{(i)} + \frac{r_k^{(i)}}{A_{kk}}e_k\right)^T
    A
    \left(e^{(i)} + \frac{r_k^{(i)}}{A_{kk}}e_k\right) \notag \\
    &=
    \Anorm{e^{(i)}}^2
    + 2\frac{r_k^{(i)}}{A_{kk}} e_k^T A e^{(i)}
    + \left(\frac{r_k^{(i)}}{A_{kk}}\right)^2 A_{kk} \notag \\
    &=
    \Anorm{e^{(i)}}^2
    - 2\frac{(r_k^{(i)})^2}{A_{kk}}
    + \frac{(r_k^{(i)})^2}{A_{kk}}.
\end{align}
Therefore
\begin{equation}
\label{eq:progress-identity}
    \Anorm{e^{(i+1)}}^2
    =
    \Anorm{e^{(i)}}^2
    -
    \frac{(r_k^{(i)})^2}{A_{kk}} .
\end{equation}
This identity is deterministic and does not depend on how the component $k$ is
chosen.  Every exact coordinate relaxation weakly decreases the $A$-norm error,
and the amount of progress is precisely the squared residual in the selected
component, scaled by the corresponding diagonal entry.

Randomization enters only through the choice of $k$.  Let
\[
    p_j^{(i)} = P(k=j \mid x^{(i)}), \qquad j=1,\ldots,n,
\]
where the distribution may be fixed in advance or may depend on the current
iterate through the residual $r^{(i)}$.  Taking conditional expectation in
\eqref{eq:progress-identity} gives the basic expected-progress formula
\begin{equation}
\label{eq:expected-progress-general}
    \E\!\left[\Anorm{e^{(i+1)}}^2 \mid x^{(i)}\right]
    =
    \Anorm{e^{(i)}}^2
    -
    \sum_{j=1}^n p_j^{(i)}
    \frac{(r_j^{(i)})^2}{A_{jj}} .
\end{equation}
Thus the analysis of any randomized component-selection rule reduces to lower
bounding the expected progress term
\[
    \sum_{j=1}^n p_j^{(i)}
    \frac{(r_j^{(i)})^2}{A_{jj}}.
\]
Uniform sampling, diagonal-weighted sampling, greedy sampling, and the
residual-weighted methods studied later differ only in how they choose the
probabilities $p_j^{(i)}$.
Note that if $r^{(i)}=0$, then $x^{(i)}=x^*$ and the method has already converged.  In
that case the sampling distribution is irrelevant.  Throughout the analysis of
residual-dependent probabilities, we implicitly restrict attention to iterations
with $r^{(i)} \neq 0$.


%\subsection{Baseline Convergence Results}

We first recall two fixed-distribution baselines.  These are the reference
points against which residual-weighted sampling will be compared.  The first
uses diagonal-weighted sampling and recovers the classical
Leventhal--Lewis-type rate for randomized coordinate descent on SPD systems.  The
second uses uniform sampling and gives the natural baseline for asynchronous
Jacobi, where uniform selection is often the simplest distribution to implement
and analyze.

\begin{theorem}[Diagonal-weighted sampling; Leventhal--Lewis]
\label{thm:leventhal-lewis}
Suppose that at each iteration,
\[
    P(k=j) = \frac{A_{jj}}{\Tr(A)}.
\]
Then
\[
    \E\!\left[\Anorm{e^{(i+1)}}^2 \mid x^{(i)}\right]
    \leq
    \left(1 - \frac{\lmin(A)}{\Tr(A)}\right)
    \Anorm{e^{(i)}}^2.
\]
\end{theorem}

\begin{theorem}[Uniform sampling]
\label{thm:uniform}
Suppose that at each iteration,
\[
    P(k=j) = \frac{1}{n}.
\]
Then
\[
    \E\!\left[\Anorm{e^{(i+1)}}^2 \mid x^{(i)}\right]
    \leq
    \left(1 - \frac{\lmin(A)}{n A_{\max}}\right)
    \Anorm{e^{(i)}}^2.
\]
\end{theorem}

Note that since $\Tr(A) = \sum_{j=1}^n A_{jj} \leq n A_{\max}$,
the diagonal-weighted bound is at least as strong as the crude uniform bound.
When $A$ has constant diagonal, the two sampling distributions coincide and
$\Tr(A)=nA_{\max}$, so the two bounds are identical.  This is the case for many
model discretizations, including standard finite-difference Laplacians with
constant diagonal.

Both theorems use the same proof template: start from the exact one-step
identity, take expectation over the selected component, and lower bound the
expected residual contribution by a multiple of $\norm{r^{(i)}}{2}^2$.  The
residual-weighted analysis below follows the same template.  The key difference
is that the sampling probabilities will depend on the current residual, causing
higher-residual components to contribute more heavily to the expected progress.
The worst-case question is whether this bias can ever hurt relative to uniform
sampling; the sharpened question is how much it helps when the residual is
non-uniform.

\subsection{Component-Selection Adaptivity Spectrum}
\label{sect:sampling-adaptivity-spectrum}

Single-component relaxation methods differ not only in the update formula, but
also in the rule used to choose the next component, row, block, or sketch.  This
choice ranges from completely oblivious rules, whose probabilities are fixed
before the iteration begins, to fully greedy rules, which compute the current
residual and select the component that promises the largest immediate progress.
The methods considered in this paper sit between these extremes: they use the
current residual to bias a randomized selection rule, but avoid the explicit
global maximization or sorting step associated with classical greedy relaxation.

At the non-adaptive end of the spectrum, the sampling distribution is fixed
before the iteration begins and never reacts to the current residual.
Strohmer and Vershynin's randomized Kaczmarz method samples rows with
probabilities proportional to squared row norms~\cite{strohmer2009randomized},
while Leventhal and Lewis analyze randomized coordinate descent and
Gauss--Seidel-type methods with probabilities chosen from static problem
data~\cite{leventhal2010randomized}. Related importance-sampling ideas appear
throughout randomized coordinate descent~\cite{nesterov2012efficiency} and
the sketch-and-project framework~\cite{gower2015randomized}. These methods
are adaptive only in the weak sense that their distributions may reflect
matrix geometry. The randomized asynchronous Jacobi analysis of
Avron, Druinsky, and Gupta~\cite{avron2015revisiting} also belongs to this
class, using uniform random selection to obtain convergence bounds under
bounded-delay and independence assumptions.

At the opposite end of the spectrum are greedy relaxation methods, which
inspect the current residual and select the component offering the largest
immediate progress. The classical Southwell or Gauss--Southwell
rule~\cite{southwell1946relaxation} and Motzkin's row-action
relaxation~\cite{agmon1954relaxation,motzkin1954relaxation} are the
historical templates; Nutini et al.~\cite{nutini2015coordinate} give a
modern coordinate-descent analysis showing that Gauss--Southwell-type
rules can converge faster than random selection except in limiting cases,
and related analyses appear in~\cite{griebel2012greedy,frommer2023convergence}.
Greedy rules are attractive because they directly target the most productive
component, but they require global residual comparisons or rank maintenance,
which can be expensive and especially awkward in parallel or asynchronous
settings.

Between these extremes lie hybrid methods that use partial residual
information. Sample-then-greedy methods like the Sampling Kaczmarz--Motzkin
algorithm of De Loera et al.~\cite{deloera2017sampling} (sharpened
by Haddock and Ma~\cite{haddock2021greed}) sample a random subset and apply
greedy selection within it. Threshold-based methods such as Bai and Wu's
greedy randomized Kaczmarz~\cite{bai2018greedyk,bai2018relaxed} and the
analogous greedy randomized coordinate descent for least squares
problems~\cite{bai2019greedy} use the residual to form an eligible set
and then sample within that set using residual-dependent weights. The
adaptive sketch-and-project framework of Gower et
al.~\cite{gower2021adaptive} provides a unifying perspective on adaptive
sampling, including max-distance, proportional, and capped sampling
variants; these correspond roughly to greedy, residual-weighted, and
thresholded selection in our taxonomy. Finally, Parallel and Distributed
Southwell~\cite{wolfsonpou2016reducing,wolfsonpou2017distributed} adapt
greedy selection for communication reduction by localizing the maximality
test. Each of these methods either uses fixed probabilities, performs
discrete greedy selection over a (possibly residual-determined) eligible
set, or maintains residual rankings; none uses a continuous residual-biased
probability distribution over all components.

The present work occupies a different point on the spectrum.  We consider
selection probabilities of the form
\[
    P(k=j) \propto |r_j|^\ell.
\]
This rule is residual-adaptive but does not require sorting, thresholding,
or constructing an eligible set.  The parameter
$\ell$ gives an explicit continuous interpolation: as $\ell \to 0$ the method
approaches uniform random selection, while as $\ell \to \infty$ it approaches
Gauss--Southwell selection. Our earlier work on residual-biased randomized
asynchronous solvers~\cite{coleman2019enhancing} is the closest predecessor:
it demonstrated empirically that non-uniform, residual-informed component
selection can improve asynchronous linear solver performance, but it used
rank-based distributions that require periodic sorting or rank maintenance.
The present work replaces that discrete rank-based mechanism with a
continuous residual-weighted sampling family and supplies a convergence
analysis explaining when such bias is beneficial.

\begin{center}
\footnotesize
\renewcommand{\arraystretch}{1.50}
\begin{tabular}{@{}p{0.20\linewidth}p{0.3\linewidth}p{0.40\linewidth}@{}}
\toprule
\textbf{Selection rule} & \textbf{Representative methods} & \textbf{Residual adaptivity} \\
\midrule
Fixed sampling
&
Randomized Kaczmarz \cite{strohmer2009randomized};
randomized CD
\cite{leventhal2010randomized,nesterov2012efficiency};
uniform asynchronous Jacobi \cite{avron2015revisiting}
&
Probabilities are fixed before the iteration begins.
The current residual does not affect the sampling rule.
\\

Sample then greed
&
Sampling Kaczmarz--Motzkin
\cite{deloera2017sampling};
improved SKM analysis \cite{haddock2021greed}
&
A random subset is sampled first; the final choice is greedy within that subset.
\\

Thresholded adaptive sampling
&
Greedy randomized Kaczmarz and CD \cite{bai2018greedyk,bai2018relaxed,bai2019greedy};
capped sketch-and-project sampling \cite{gower2021adaptive}
&
The residual determines an eligible set or capped set, and sampling is then
performed within that set using residual-dependent weights.
\\

Greedy selection
&
Southwell relaxation \cite{southwell1946relaxation};
Motzkin relaxation \cite{agmon1954relaxation,motzkin1954relaxation};
Gauss--Southwell CD \cite{nutini2015coordinate};
Rule-based greedy approach \cite{griebel2012greedy,frommer2023convergence}
&
The selected component maximizes a residual-based progress measure.
\\

Parallel greedy selection
&
Parallel Southwell and Distributed Southwell
\cite{wolfsonpou2016reducing,wolfsonpou2017distributed}
&
Greedy selection is localized to reduce communication.
\\

Rank-based residual-biased sampling
&
Residual-biased randomized asynchronous solvers
\cite{coleman2019enhancing}
&
Components ranked by current residual magnitude and a non-uniform distribution is used to sample.
\\

Continuous residual-weighted sampling
&
This work; adaptive sketch-and-project proportional sampling
\cite{gower2021adaptive}
&
All components eligible, probabilities vary directly \& smoothly with current residual.  Parameters give explicit limits from uniform
sampling to Gauss--Southwell selection.
\\
\bottomrule
\end{tabular}
\end{center}


\section{Residual-Weighted Sampling}
\label{sect:baseline-residual-weighted}

We establish that residual-weighted sampling matches the uniform bound in the worst case. The following two theorems are, to our knowledge, not stated explicitly in the prior literature; they are recorded here because they serve as the starting point for the sharpened bounds in Section~\ref{sec:sharp}. Uniform sampling emerges as a special case of Theorem~\ref{thm:baseline-power} as $\ell \to 0$.

\begin{theorem}[Baseline bound, power-weighted sampling]
\label{thm:baseline-power}
Under power-weighted sampling $P(k=j) = |r_j^{(i)}|^\ell / \sum_m |r_m^{(i)}|^\ell$
with any $\ell > 0$:
$$
\E[\Anorm{e^{(i+1)}}^2 \mid x^{(i)}] \leq \left(1 - \frac{\lmin(A)}{n \cdot A_{\max}}\right) \Anorm{e^{(i)}}^2.
$$
\end{theorem}

\begin{proof}
From~\eqref{eq:expected-progress-general}, the expected progress is
$\E[\mathrm{Progress}] = \sum_j P(j)\,(r_j^{(i)})^2/A_{jj}$.
Bounding $A_{jj} \leq A_{\max}$:
$$
\E[\mathrm{Progress}] \;\geq\;
\frac{1}{A_{\max}}\cdot\frac{\sum_j |r_j^{(i)}|^{\ell+2}}{\sum_m |r_m^{(i)}|^\ell}
\;=\;
\frac{1}{A_{\max}}\cdot\frac{\|r^{(i)}\|_{\ell+2}^{\ell+2}}{\|r^{(i)}\|_\ell^\ell}.
$$
It remains to show that
\begin{equation}
\label{eq:norm-ratio-bound}
\frac{\|r\|_{\ell+2}^{\ell+2}}{\|r\|_\ell^\ell} \;\geq\; \frac{\|r\|_2^2}{n}
\qquad \text{for every nonzero } r \in \R^n \text{ and every } \ell > 0.
\end{equation}
Define $p_j = r_j^2/\|r\|_2^2$, which is a probability distribution on $\{1,\ldots,n\}$.
Then $|r_j|^{2k} = p_j^k\,\|r\|_2^{2k}$ for every $k \geq 0$, so
\[
\|r\|_\ell^\ell = \|r\|_2^\ell \sum_j p_j^{\ell/2},
\qquad
\|r\|_{\ell+2}^{\ell+2} = \|r\|_2^{\ell+2} \sum_j p_j^{(\ell+2)/2},
\]
and~\eqref{eq:norm-ratio-bound} reduces to
\begin{equation}
\label{eq:cheb-target}
n \sum_j p_j^{(\ell+2)/2} \;\geq\; \sum_j p_j^{\ell/2},
\end{equation}
since substituting the norm identities into~\eqref{eq:norm-ratio-bound}
and simplifying by $\|r\|_2^2$ gives
$\sum_j p_j^{(\ell+2)/2} / \sum_j p_j^{\ell/2} \geq 1/n$.
Setting $\alpha = \ell/2 \geq 0$, both sequences $(p_j)_{j=1}^n$ and $(p_j^\alpha)_{j=1}^n$ are
similarly ordered (the function $t \mapsto t^\alpha$ is non-decreasing on $[0,1]$ for $\alpha \geq 0$).
Chebyshev's sum inequality\footnote{For real sequences
$a_1 \leq \cdots \leq a_n$ and $b_1 \leq \cdots \leq b_n$,
$n\sum_i a_i b_i \geq (\sum_i a_i)(\sum_i b_i)$; see e.g.
\cite[Theorem 43]{hardy1952inequalities}.} therefore gives
\[
n \sum_j p_j \cdot p_j^\alpha
\;\geq\;
\Bigl(\sum_j p_j\Bigr)\Bigl(\sum_j p_j^\alpha\Bigr)
\;=\; \sum_j p_j^\alpha,
\]
using $\sum_j p_j = 1$. This is~\eqref{eq:cheb-target}, hence~\eqref{eq:norm-ratio-bound}.
Combining with $\|r^{(i)}\|_2^2 \geq \lmin(A)\,\Anorm{e^{(i)}}^2$ yields the claim.
\end{proof}

Theorem~\ref{thm:baseline-power} gives the \emph{same} worst-case bound as uniform sampling (Theorem~\ref{thm:uniform}). Residual-weighting is guaranteed not to hurt, but these baseline bounds provide no incentive to use it. Section~\ref{sec:sharp} shows that the gap between these bounds and the empirical convergence rate is quantified by the inverse participation ratio.


\subsection{Sharpened Bounds via the Inverse Participation Ratio}
\label{sec:sharp}

\begin{definition}[Inverse Participation Ratio]
\label{def:ipr}
For a nonzero vector $r \in \R^n$, the \emph{inverse participation ratio} (IPR) is
\begin{equation}
\label{eq:ipr}
\nu^2(r) = \frac{n \cdot \norm{r}{4}^4}{\norm{r}{2}^4} = \frac{n \sum_j r_j^4}{\left(\sum_j r_j^2\right)^2}.
\end{equation}
\end{definition}

The squared form reflects the fact that the quantity naturally enters bounds
as $\nu^2 = (\sqrt{n}\,\|r\|_4/\|r\|_2)^2$; writing it unsquared as
$\nu = \sqrt{n}\,\|r\|_4/\|r\|_2$ is occasionally useful (e.g., in
Cauchy--Schwarz bounds) but we use $\nu^2$ by default. The IPR has the
following properties.

\begin{lemma}[Properties of the IPR]
\label{lem:ipr-properties}
For any nonzero $r \in \R^n$:
\begin{enumerate}[label=(\roman*)]
    \item $\nu^2(r) \geq 1$, with equality if and only if all $|r_j|$ are equal.
    \item $\nu^2(r) \leq n$, with equality if and only if $r$ has exactly one nonzero component.
    \item $\nu^2(r)$ is invariant under uniform scaling: $\nu^2(cr) = \nu^2(r)$ for $c \neq 0$.
    \item $\nu^2(r)$ depends only on the \emph{shape} of the distribution of $|r_j|^2$, not its magnitude.
\end{enumerate}
\end{lemma}

\begin{proof}
\textit{(i)} By Cauchy--Schwarz applied to the vectors
$(1, \ldots, 1)$ and $(r_1^2, \ldots, r_n^2)$,
\[
\Bigl(\sum_j r_j^2\Bigr)^2 = \Bigl(\sum_j 1 \cdot r_j^2\Bigr)^2
\;\leq\; n \sum_j r_j^4,
\]
giving $\nu^2(r) \geq 1$. Equality in Cauchy--Schwarz holds iff
$(r_1^2, \ldots, r_n^2)$ is a scalar multiple of $(1, \ldots, 1)$,
i.e., iff all $|r_j|$ are equal.

\textit{(ii)} For any $r \in \R^n$,
\[
\sum_j r_j^4
\;\leq\; (\max_j r_j^2) \sum_j r_j^2
\;\leq\; \Bigl(\sum_j r_j^2\Bigr)^2,
\]
where the first inequality bounds each $r_j^2$ in $r_j^4 = r_j^2 \cdot r_j^2$
by the maximum, and the second uses $\max_j r_j^2 \leq \sum_j r_j^2$ for
non-negative terms. Hence $\nu^2(r) \leq n$. Equality throughout requires
$\max_j r_j^2 = \sum_j r_j^2$, which (for non-negative summands) holds iff
exactly one $r_j$ is nonzero.

\textit{(iii)} For $c \neq 0$,
$\nu^2(cr) = n\,c^4 \|r\|_4^4 / (c^2 \|r\|_2^2)^2 = \nu^2(r)$.

\textit{(iv)} Follows from (iii): the normalized squared residual
$p_j = r_j^2/\|r\|_2^2$ is invariant under uniform scaling, and a direct
computation gives $\nu^2(r) = n\sum_j p_j^2$, which depends only on $(p_j)$.
\end{proof}

The particular form of $\nu^2$ is not a design choice but a consequence of
the proof. The one-step progress identity~\eqref{eq:progress-identity}
combined with power-weighted sampling $P(k=j) = r_j^2/\norm{r}{2}^2$
produces the ratio $\norm{r}{4}^4/\norm{r}{2}^2$ as a natural progress
lower bound (derivation in the proof of Theorem~\ref{thm:sharp-power}).
The question of what this ratio equals is forced: to compare with uniform
sampling we need to express it as a multiple of $\norm{r}{2}^2/n$, and
the multiple that appears is exactly $\nu^2$.
Other scalar measures of non-uniformity are possible, but this one is forced by
the expected-progress identity; attempts to use
$\|r\|_\infty / (\|r\|_2/\sqrt n)$ or $\|r\|_1 / (\sqrt n \|r\|_2)$
produce measures of non-uniformity but do not arise from the proof
and do not give the sharpened bound.

The IPR admits three complementary interpretations.
As a \emph{localization or participation} measure, it is a standard tool
in physics for quantifying how many sites or degrees of freedom meaningfully
participate in a state~\cite{evers2008anderson}: $\nu^2 \approx 1$ means
delocalized, $\nu^2 \gg 1$ means localized on a few components.
As an \emph{effective support}, treating $p_j = r_j^2/\|r\|_2^2$ as a
probability distribution gives $\nu^2(r) = n\sum_j p_j^2$, so
$n/\nu^2(r) = 1/\sum_j p_j^2 = e^{H_2(p)}$ is the order-2 effective
support of $p$, where $H_2(p) = -\log\sum_j p_j^2$ is the R\'enyi-2
entropy~\cite{renyi1961measures}; a residual spread uniformly over $5$
components has $\nu^2 = n/5$.
Finally, a direct computation gives $\nu^2(r) = 1 + \mathrm{Var}(r_j^2)/\bar{r^2}^{\,2}$
with mean and variance taken over a uniform index $j$, so $\nu^2$ is the
\emph{squared coefficient of variation} of $(r_j^2)$ and $\nu^2 > 1$ iff
the squared residuals are not all equal. For our purposes:
$\nu^2 \approx 1$ means biased sampling has little advantage, while
$\nu^2 \gg 1$ means a few components dominate and biased sampling can
concentrate effort where it matters. 
% The steady-state IPR values on our
% three test problems (full details in Section~\ref{sec:experiments})
% are $\nu^2 \approx 4.8$ for \texttt{laplace}, $\nu^2 \approx 5.7$ for
% \texttt{poisson}, and $\nu^2 \approx 4.0$ for \texttt{fem}, corresponding
% to effective supports of a few thousand grid points out of tens of thousands.


\subsection{Sharpened Convergence for Power-Weighted Jacobi ($\ell = 2$)}

While the baseline bounds in the previous section apply to multiple weighting parameters, the choice $\ell=2$ possesses a unique analytical advantage. When sampling probabilities are proportional to the squared residual components, $P(k=j) = r_j^2/\|r\|_2^2$, they form a distribution whose expected progress naturally scales with the variance of those squared residuals. This choice causes the non-uniformity of the residual to drop directly out of the one-step expectation as a clean scalar multiplier, allowing the inverse participation ratio to act as the exact bridge between the empirical spread of the residual and the theoretical convergence rate.

\begin{theorem}[Sharpened Power-Weighted Convergence]
\label{thm:sharp-power}
Let $A \in \R^{n \times n}$ be SPD. Under power-weighted randomized Jacobi with $\ell = 2$ (probability $\propto r_j^2$):
\begin{equation}
\label{eq:sharp-power}
\E[\Anorm{e^{(i+1)}}^2 \mid x^{(i)}] \leq \left(1 - \frac{\nu_i^2 \cdot \lmin(A)}{n \cdot A_{\max}}\right) \Anorm{e^{(i)}}^2
\end{equation}
where $\nu_i^2 = \nu^2(r^{(i)}) = n\norm{r^{(i)}}{4}^4 / \norm{r^{(i)}}{2}^4$ is the IPR of the current residual.

In particular, the expected per-step progress is amplified by $\nu_i^2$ relative to uniform sampling, replacing the uniform contraction factor of $(1-\alpha)$ by $(1 - \nu_i^2\alpha)$ where $\alpha = \lmin(A)/nA_{\max}$.
\end{theorem}

\begin{proof}
Write $\nu^2 = \nu^2(r^{(i)})$ throughout. From~\eqref{eq:expected-progress-general},
the expected progress under $P(k=j) = r_j^2/\|r\|_2^2$ is
\begin{align}
\E[\mathrm{Progress} \mid x^{(i)}]
&= \sum_{j=1}^n \frac{r_j^2}{\|r\|_2^2}\cdot \frac{r_j^2}{A_{jj}}
 = \frac{1}{\|r\|_2^2}\sum_{j=1}^n \frac{r_j^4}{A_{jj}}.
\end{align}
Using $A_{jj} \leq A_{\max}$:
\begin{equation}
\E[\mathrm{Progress} \mid x^{(i)}]
\;\geq\; \frac{1}{A_{\max}}\cdot \frac{\|r\|_4^4}{\|r\|_2^2}.
\end{equation}
By the IPR identity $\|r\|_4^4 = (\nu^2/n)\,\|r\|_2^4$
(Definition~\ref{def:ipr}):
\begin{equation}
\E[\mathrm{Progress} \mid x^{(i)}]
\;\geq\; \frac{\nu^2}{n A_{\max}}\,\|r\|_2^2.
\end{equation}
Applying $\|r\|_2^2 \geq \lmin(A)\Anorm{e^{(i)}}^2$ from~\eqref{eq:residual-energy-bounds}:
\[
\E[\mathrm{Progress} \mid x^{(i)}]
\;\geq\; \frac{\nu^2\,\lmin(A)}{n\,A_{\max}}\,\Anorm{e^{(i)}}^2,
\]
and the claim follows from $\E[\Anorm{e^{(i+1)}}^2 \mid x^{(i)}] =
\Anorm{e^{(i)}}^2 - \E[\mathrm{Progress} \mid x^{(i)}]$.
\end{proof}

Setting $\nu_i^2 = 1$ recovers the baseline bound of
Theorem~\ref{thm:baseline-power} exactly, so the sharpened bound is
never weaker than the baseline and is strictly stronger whenever
$\nu_i > 1$. In an implementation, $\nu_i^2$ costs one additional pass
over the residual to compute, on top of the pass already needed for
the sampling step itself, so the IPR doubles as a per-iteration
diagnostic that reports in real time how much faster the iteration is
running than the uniform rate would predict. For concreteness, on the
\texttt{poisson} problem of Section~\ref{sec:experiments}
($n = 16{,}384$, $\lambda_{\min}/A_{\max} \approx 10^{-3}$) the
steady-state IPR $\nu^2 \approx 5.7$ amplifies the uniform-sampling
progress coefficient by roughly a factor of six, and the
early-iteration transient (where $\nu^2$ is in the tens to hundreds)
contributes the bulk of the actual residual reduction;
Figure~\ref{fig:combined_sync_bounds} shows that the realized
convergence on this problem is in fact faster than the steady-state
bound alone predicts, precisely because of this transient.


%\subsection{Global Convergence Under a Sustained Non-Uniformity Assumption}

The one-step result (Theorem~\ref{thm:sharp-power}) holds unconditionally. To obtain a global (multi-step) rate, we need the IPR to remain bounded below across iterations.

\begin{corollary}[Global Rate Under Sustained Non-Uniformity]
\label{cor:global-rate}
If $\nu_i^2 \geq \bar{\nu}^2 > 1$ for all $i = 0, 1, \ldots, m-1$, then
$$
\E[\Anorm{e^{(m)}}^2] \leq \left(1 - \frac{\bar{\nu}^2 \cdot \lmin(A)}{n \cdot A_{\max}}\right)^m \Anorm{e^{(0)}}^2.
$$
\end{corollary}

Whether $\nu_i^2 \geq \bar{\nu}^2$ holds in practice depends on the problem. We investigate this numerically in Section~\ref{sec:experiments}. Proving it analytically for specific matrix classes (e.g., discretized PDEs with localized forcing) is an interesting open direction.
Without the sustained non-uniformity assumption, we can still obtain a global bound by using the IPR as a per-step weight.

\begin{corollary}[Global Rate via Geometric Mean of IPR]
\label{cor:geometric-mean}
After $m$ iterations,
$$
\E[\Anorm{e^{(m)}}^2] \leq \prod_{i=0}^{m-1}\left(1 - \frac{\nu_i^2 \cdot \lmin(A)}{n \cdot A_{\max}}\right) \Anorm{e^{(0)}}^2.
$$
\end{corollary}

\noindent
Taking logarithms, the effective convergence rate is governed by the arithmetic mean of $\log(1 - \nu_i^2 \lmin/(nA_{\max}))$, which is dominated by the average IPR.


\section{Asynchronous Power-Weighted Jacobi}
\label{sec:async}

The synchronous result in Theorem~\ref{thm:sharp-power} sharpens the
uniform-sampling bound by exactly the IPR, giving a clean theoretical
separation between residual-weighted and uniform sampling at every
step. A natural next question is whether the same analytical leverage
carries over to asynchronous execution, where the residual driving the
sampling distribution may be stale relative to the iterate being
updated. The question is not merely academic: residual-weighted
sampling is most appealing precisely in the high-concurrency regime
where asynchrony is unavoidable, and the synchronous bound is silent
about what happens there.

This section extends the sharpened synchronous analysis to the
asynchronous setting. We first review the Avron--Druinsky--Gupta (ADG)
framework~\cite{avron2015revisiting} for uniform asynchronous Jacobi,
then identify three specific obstacles that arise under
residual-weighted sampling and show how each can be handled under a
global bound on the IPR trajectory. The result is a convergence
theorem structurally analogous to ADG's Theorem~2a, with the
uniform-sampling cross-term parameter $\rho$ replaced by a dynamic
parameter $\rho_{\text{dyn}}$ and with the IPR entering both the
progress coefficient and the allowed-delay window. The theorem should
be read as a first asynchronous convergence result for power-weighted
sampling rather than as a strict generalization of the uniform theory;
several aspects of the bound are provably loose and are flagged as
sharpening targets at the end of the section.

%\subsection{The Asynchronous Iteration}

To match the notation of Section~\ref{sec:background}, we write the asynchronous iteration in component form. At step $j$ the worker reads a stale iterate $x^{(s(j))}$ (where $s(j) \leq j$ is the step at which the read completed), samples a component $k_j$ according to a probability distribution that may depend on the stale residual $\tilde r^{(j)} = b - A x^{(s(j))}$, computes the Jacobi update using the stale residual's $k_j$-th component, and commits the result to the current iterate $x^{(j)}$:
\begin{equation}
\label{eq:async-update}
x^{(j+1)} = x^{(j)} + \beta\cdot\frac{\tilde r_{k_j}^{(j)}}{A_{k_j, k_j}}\, e_{k_j}.
\end{equation}
The bounded-delay condition $j - s(j) \leq \tau$ limits how stale any read can be; $\beta \in (0, 1]$ is an optional step-size parameter, with $\beta = 1$ corresponding to a full Jacobi step. Setting $s(j) = j$ for all $j$ (no staleness) and $\beta = 1$ recovers the synchronous iteration~\eqref{eq:update}. The update is applied to $x^{(j)}$, not $x^{(s(j))}$; the read is stale but the write is current.

\begin{algorithm}[ht]
\small
\caption{Asynchronous power-weighted randomized Jacobi (per-worker loop)}
\label{alg:async-power-jacobi}
\begin{algorithmic}[1]
\Require SPD matrix $A$, RHS $b$, shared iterate $x$, shared
         residual $r$, exponent $\ell > 0$, step size
         $\beta \in (0, 1]$, read mode $\in$ \{consistent,
         inconsistent\}
\While{not converged}
  \If{read mode is consistent}
    \State $r_{\text{local}} \gets$ atomic snapshot of $r$
    \Comment{Used for both sampling and update value}
  \Else
    \State $r_{\text{local}} \gets r$ via \texttt{\#pragma omp atomic read}
    \Comment{Live; second read happens after sampling}
  \EndIf
  \State Sample $k \in \{1, \ldots, n\}$ with
         $P(k = j) = |r_{\text{local},j}|^\ell / \sum_m |r_{\text{local},m}|^\ell$
  \If{read mode is inconsistent}
    \State $r_k \gets$ atomic read of $r_k$
    \Comment{Second, possibly-stale-different read}
  \Else
    \State $r_k \gets r_{\text{local},k}$
  \EndIf
  \State $\delta \gets \beta \cdot r_k / A_{kk}$
  \State Atomic update: $x_k \gets x_k + \delta$
  \State Atomically apply stencil updates to $r$ at $k$ and
         its neighbors in the matrix graph
\EndWhile
\end{algorithmic}
\end{algorithm}

For comparison with the ADG literature, our $x^{(s(j))}$ corresponds to their $x_{k(j)}$ and our $k_j$ to their direction $d_j = e_{k_j}$; our $\tilde r_{k_j}^{(j)}/A_{k_j,k_j}$ corresponds to their $\gamma_j = (x^* - x_{k(j)}, d_j)_A$ (for unit-diagonal $A$ and $d_j = e_{k_j}$, $(x^* - x_{k(j)}, e_{k_j})_A = [A(x^* - x_{k(j)})]_{k_j} = \tilde r_{k_j}^{(j)}$). Translation between the two conventions is purely notational.
%\subsection{Review of the Avron--Druinsky--Gupta Framework}
Under uniform direction sampling ($P(k_j = r) = 1/n$), consistent reads, and an independent-delays assumption (ADG Assumption A-4), Avron et al.\ establish the following result.

\begin{theorem}[ADG Theorem~2a]
\label{thm:avron-async}
Consider iteration~\eqref{eq:async-update} with uniform sampling, $\beta = 1$, unit diagonal $A$, consistent reads, and bounded delay $\tau$. Let $\rho = \frac{1}{n}\|A\|_\infty$. Provided $2\rho\tau < 1$, for every epoch of length $m \geq 0.693\,n/\lmax$,
$$
E_m \leq \left(1 - \frac{1 - 2\rho\tau}{2\kappa}\right) E_0,
$$
where $E_j = \E[\Anorm{e^{(j)}}^2]$.
\end{theorem}

The proof architecture has three stages: a one-step progress lemma with
both lower and upper bounds on the expected squared inner product at
the stale state; an unrolled recursion that separates per-step
progress from cross-coupling damage between stale components; and an
epoch assembly that sums the recursion using the upper bound to ensure
the error does not collapse faster than the damage can accumulate.

%\subsection{Three Obstacles Under Residual-Weighted Sampling}

Porting Theorem~\ref{thm:sharp-power} ($\ell = 2$) into this architecture, three specific mechanisms break.

\paragraph{Obstacle 1: The independent-delays assumption fails.} ADG Assumption A-4 requires that $s(j)$ is independent of the sequence of component choices $k_0, k_1, \ldots$. Under uniform sampling this is plausible; the scheduler's choice of when to read cannot influence which component gets sampled. Under power-weighted sampling the distribution over components at step $j$ is $P(k_j = r) = (\tilde r_r^{(j)})^2 / \|\tilde r^{(j)}\|_2^2$, which depends on the stale residual and hence on $s(j)$. The component choice and the delay are coupled through the stale iterate.

\paragraph{Obstacle 2: The cross-term bound loses its $1/n$ scaling.} In the uniform proof, for any fixed past component $k_t = l$,
$$
\E[\,|A_{l, k_j}|\,] = \frac{1}{n}\sum_r |A_{lr}| \leq \frac{1}{n}\|A\|_\infty = \rho.
$$
The $1/n$ factor is what makes $\rho = O(1/n)$ for sparse matrices. Under power-weighted sampling the analogous expectation is $\sum_r ((\tilde r_r^{(j)})^2/\|\tilde r^{(j)}\|_2^2)\, |A_{lr}|$ and the favorable $1/n$ prefactor is lost. A residual concentrated on a single component would put all its probability mass on a single matrix row, giving no $n$-dependence at all.

\paragraph{Obstacle 3: The progress lemma needs an upper bound.} Theorem~\ref{thm:sharp-power} provides a \emph{lower} bound on expected progress. The async epoch assembly also requires an \emph{upper} bound to guarantee that delayed error terms are still large enough fractions of the initial error for the sum over the epoch to accumulate.

%\subsection{Handling the Obstacles via a Global IPR Bound}
\vspace{1em}

Our main tool for handling these obstacles is to assume a uniform bound on the IPR throughout the epoch. This simultaneously (a) assumes away Obstacle 1 by imposing worst-case dynamics rather than tracking actual correlations, (b) partially repairs Obstacle 2 by allowing Cauchy--Schwarz to decouple residual from matrix, and (c) gives us a two-sided lemma that resolves Obstacle 3 directly.

\begin{definition}[Uniform IPR trajectory]
\label{assump:global-ipr}
There exist constants $1 \leq \nu_{\min}^2 \leq \nu_{\max}^2 \leq n$ such that for every step $j$ in the epoch,
$$
\nu_{\min}^2 \leq \nu^2(\tilde r^{(j)}) \leq \nu_{\max}^2.
$$
\end{definition}

Definition~\ref{assump:global-ipr} requires a priori bounds on the full
trajectory of IPRs, including during early iterations when the
residual non-uniformity is actively growing (see
Section~\ref{sec:experiments}). In an implementation one could compute
$\nu^2(\tilde r^{(j)})$ online and verify the assumption post hoc, but
it is an input to the theorem rather than an output.
Section~\ref{sec:experiments} verifies the assumption empirically
across all three test problems, with sustained $\nu^2$ values that
vary by less than $3\%$ across two decades of thread count
(Table~\ref{tab:hw-scaling}). What we get in exchange is a mechanical
path to a convergence theorem with the right structural form; closing
the gap between the assumption and the dynamics analytically is the
main open problem.

With $\ell = 2$ power-weighting at the stale state, letting $\tilde e^{(j)} = x^{(s(j))} - x^*$,
$$
\E[(\tilde e^{(j)}, e_{k_j})_A^2] = \sum_{r=1}^n \frac{(\tilde r_r^{(j)})^2}{\|\tilde r^{(j)}\|_2^2}\, (\tilde r_r^{(j)})^2 = \frac{\|\tilde r^{(j)}\|_4^4}{\|\tilde r^{(j)}\|_2^2} = \frac{\nu^2(\tilde r^{(j)})}{n}\|\tilde r^{(j)}\|_2^2.
$$
Applying $\lmin \Anorm{\tilde e^{(j)}}^2 \leq \|\tilde r^{(j)}\|_2^2 \leq \lmax \Anorm{\tilde e^{(j)}}^2$ gives a two-sided control on the conditional expected progress.

\begin{lemma}[Dynamic one-step lemma]
\label{lem:dynamic-async-1}
Under power-weighted sampling with $\ell = 2$,
$$
\frac{\nu^2(\tilde r^{(j)})\,\lmin}{n}\Anorm{\tilde e^{(j)}}^2 \;\leq\; \E[(\tilde e^{(j)}, e_{k_j})_A^2] \;\leq\; \frac{\nu^2(\tilde r^{(j)})\,\lmax}{n}\Anorm{\tilde e^{(j)}}^2.
$$
Under Definition~\ref{assump:global-ipr}, the same bounds hold with $\nu^2(\tilde r^{(j)})$ replaced by $\nu_{\min}^2$ (lower) and $\nu_{\max}^2$ (upper).
\end{lemma}

For the cross-term, under power-weighted sampling, for past component $k_t = l$ fixed,
$$
\E[\,|A_{l, k_j}|\,] = \sum_r \frac{(\tilde r_r^{(j)})^2}{\|\tilde r^{(j)}\|_2^2}\,|A_{lr}|
\leq \sqrt{\sum_r \left(\frac{(\tilde r_r^{(j)})^2}{\|\tilde r^{(j)}\|_2^2}\right)^{\!2}}\cdot\|A_{l,:}\|_2
= \frac{\nu(\tilde r^{(j)})}{\sqrt n}\,\|A_{l,:}\|_2
$$
by Cauchy--Schwarz. Defining $\rho_2 := \max_l \|A_{l,:}\|_2$ and using Definition~\ref{assump:global-ipr},
$$
\E[\,|A_{l, k_j}|\,] \leq \frac{\nu_{\max}\,\rho_2}{\sqrt n} =: \rho_{\text{dyn}}.
$$
For a sparse matrix with $c$ nonzeros per row and bounded entries, $\rho \approx cM/n = O(1/n)$ in the uniform case, while $\rho_{\text{dyn}} = \nu_{\max} M\sqrt c/\sqrt n = O(\nu_{\max}/\sqrt n)$ here. Even setting $\nu_{\max} = 1$ (flat residual) gives $\rho_{\text{dyn}} = O(1/\sqrt n)$, which is strictly worse than the ADG bound of $O(1/n)$. The dynamic analysis does not recover the uniform result as a special case; the Cauchy--Schwarz step throws away $\ell_1$-style summability that the uniform analysis uses directly. This is the most important sharpening target (see Section~\ref{sec:limitations}).


\subsection{Convergence Theorem}

\begin{theorem}[Asynchronous power-weighted Jacobi, $\ell = 2$]
\label{thm:async-power-weighted}
Consider the asynchronous iteration~\eqref{eq:async-update} with step-size $\beta = 1$, unit diagonal $A$, power-weighted sampling ($\ell = 2$, using the stale residual), consistent reads, and bounded delay $\tau$. Under Definition~\ref{assump:global-ipr}, let
$$
\rho_{\text{dyn}} = \frac{\nu_{\max}\,\rho_2}{\sqrt n}, \qquad \tilde\nu_\tau = 1 - 2\rho_{\text{dyn}}\tau.
$$
Provided $\tau < \sqrt n / (2\nu_{\max}\rho_2)$, for every epoch of length $m \geq 0.693\,n/(\nu_{\max}^2\lmax)$,
$$
E_m \leq \left(1 - \frac{\nu_{\min}^2}{\nu_{\max}^2} \cdot \frac{\tilde\nu_\tau}{2\kappa}\right) E_0.
$$
\end{theorem}

\begin{proof}
The proof adapts the three-stage architecture of ADG~\cite{avron2015revisiting} (Theorem~2) to the power-weighted setting: one-step expansion, telescoped recursion with cross-term decoupling, and epoch assembly via the ``error cannot collapse'' argument. The substitutions are: Lemma~\ref{lem:dynamic-async-1} replaces the uniform one-step lemma, and the cross-term parameter $\rho_{\text{dyn}}$ from Section~\ref{sec:async} replaces ADG's $\rho$. We carry out the steps that involve the residual-weighted analysis and refer to ADG for the standard epoch-assembly.

\smallskip\noindent\textit{Step 1: One-step error identity.} With unit-diagonal $A$ and $\beta = 1$, the update~\eqref{eq:async-update} gives $e^{(j+1)} = e^{(j)} + \tilde r_{k_j}^{(j)} e_{k_j}$. The accumulated in-flight updates between read time $s(j)$ and write time $j$ give
\[
r^{(j)} = \tilde r^{(j)} - \sum_{l \in J_j} \tilde r_{k_l}^{(l)} A e_{k_l}, \qquad J_j = \{l : s(j) \leq l < j\}, \qquad |J_j| \leq \tau,
\]
where $J_j$ is the in-flight set and the bound on its size is the bounded-delay condition. Expanding $\Anorm{e^{(j+1)}}^2$ in the $A$-inner product, using $A_{k_j,k_j} = 1$ and substituting $(Ae^{(j)})_{k_j} = -r_{k_j}^{(j)}$,
\begin{equation}
\label{eq:async-onestep-proof}
\Anorm{e^{(j+1)}}^2 = \Anorm{e^{(j)}}^2 - (\tilde r_{k_j}^{(j)})^2 + 2 \sum_{l \in J_j} A_{k_j, k_l}\,\tilde r_{k_j}^{(j)}\,\tilde r_{k_l}^{(l)}.
\end{equation}
The negative term is the stale progress; the sum is the interference from in-flight updates.

\smallskip\noindent\textit{Step 2: AM--GM decoupling.} Apply $2|ab| \leq a^2 + b^2$ to each cross-coupling term:
\begin{equation}
\label{eq:async-amgm}
\left|2 A_{k_j, k_l}\,\tilde r_{k_j}^{(j)}\,\tilde r_{k_l}^{(l)}\right| \leq |A_{k_j, k_l}| \left[(\tilde r_{k_j}^{(j)})^2 + (\tilde r_{k_l}^{(l)})^2\right].
\end{equation}

\smallskip\noindent\textit{Step 3: Conditional expectations.} Let $\mathcal{F}_j$ denote the $\sigma$-algebra of all information up to (but not including) the random choice of $k_j$, so that $\tilde r^{(j)}$, $J_j$, and $\{(k_l, \tilde r^{(l)}) : l < j\}$ are $\mathcal{F}_j$-measurable. We require three conditional bounds.

\textit{(a) Progress (lower).} Lemma~\ref{lem:dynamic-async-1} gives
\[
\E\!\left[(\tilde r_{k_j}^{(j)})^2 \mid \mathcal{F}_j\right] \;\geq\; \frac{\nu_{\min}^2\,\lmin}{n}\,\Anorm{\tilde e^{(j)}}^2.
\]

\textit{(b) Stale interference factor.} Since $(\tilde r_{k_l}^{(l)})^2$ is $\mathcal{F}_j$-measurable it pulls out of the expectation, and the matrix entry is bounded by the cross-term inequality preceding Theorem~\ref{thm:async-power-weighted}:
\[
\E\!\left[|A_{k_j, k_l}|\,(\tilde r_{k_l}^{(l)})^2 \mid \mathcal{F}_j\right] = (\tilde r_{k_l}^{(l)})^2\,\E\!\left[|A_{k_j, k_l}| \mid \mathcal{F}_j\right] \;\leq\; \rho_{\text{dyn}}\,(\tilde r_{k_l}^{(l)})^2.
\]

\textit{(c) Current interference factor.} The sampling distribution is correlated with the residual square, so the bound proceeds through the IPR identity rather than the Cauchy--Schwarz cross-term inequality alone:
\begin{align*}
\E\!\left[|A_{k_j, k_l}|\,(\tilde r_{k_j}^{(j)})^2 \mid \mathcal{F}_j\right]
&= \sum_{r=1}^n \frac{(\tilde r_r^{(j)})^2}{\|\tilde r^{(j)}\|_2^2}\,|A_{k_l, r}|\,(\tilde r_r^{(j)})^2
= \frac{1}{\|\tilde r^{(j)}\|_2^2}\sum_{r=1}^n |A_{k_l, r}|\,(\tilde r_r^{(j)})^4 \\
&\leq \|A_{k_l, :}\|_\infty \cdot \frac{\|\tilde r^{(j)}\|_4^4}{\|\tilde r^{(j)}\|_2^2}
= \|A_{k_l, :}\|_\infty \cdot \frac{\nu^2(\tilde r^{(j)})}{n}\,\|\tilde r^{(j)}\|_2^2 \\
&\leq \rho_2 \cdot \frac{\nu_{\max}^2}{n}\,\|\tilde r^{(j)}\|_2^2 \;\leq\; \rho_{\text{dyn}}\,\|\tilde r^{(j)}\|_2^2,
\end{align*}
where we used $|A_{k_l, r}| \leq \|A_{k_l, :}\|_\infty \leq \|A_{k_l, :}\|_2 \leq \rho_2$, the IPR identity $\|\tilde r^{(j)}\|_4^4 = (\nu^2(\tilde r^{(j)})/n)\,\|\tilde r^{(j)}\|_2^4$, the global bound $\nu^2 \leq \nu_{\max}^2$, and finally $\nu_{\max}^2/n \leq \nu_{\max}/\sqrt n$ (which holds because $\nu_{\max} \leq \sqrt n$).

Combining (a)--(c) with~\eqref{eq:async-onestep-proof} and~\eqref{eq:async-amgm}, and using $|J_j| \leq \tau$,
\begin{equation}
\label{eq:async-condexp}
\E[\Anorm{e^{(j+1)}}^2 \mid \mathcal{F}_j] \;\leq\; \Anorm{e^{(j)}}^2 - \frac{\nu_{\min}^2\,\lmin}{n}\,\Anorm{\tilde e^{(j)}}^2 + \rho_{\text{dyn}}\tau\,\|\tilde r^{(j)}\|_2^2 + \rho_{\text{dyn}}\sum_{l \in J_j}(\tilde r_{k_l}^{(l)})^2.
\end{equation}

\smallskip\noindent\textit{Step 4: Telescoping and double-counting.} Take total expectation, write $E_j = \E[\Anorm{e^{(j)}}^2]$, and sum~\eqref{eq:async-condexp} over $j = 0, \ldots, m-1$. The double sum over in-flight pairs $(j, l \in J_j)$ re-indexes as a single sum over $l$, since each $l$ appears in $J_j$ for at most $\tau$ subsequent steps:
\[
\sum_{j=0}^{m-1} \sum_{l \in J_j} \E[(\tilde r_{k_l}^{(l)})^2] \;\leq\; \tau \sum_{l=0}^{m-1} \E[(\tilde r_{k_l}^{(l)})^2].
\]
Both single sums in the resulting expression involve $\E[(\tilde r_{k_j}^{(j)})^2]$ and $\E[\|\tilde r^{(j)}\|_2^2]$, which are controlled by $\E[\Anorm{\tilde e^{(j)}}^2]$ via the upper side of Lemma~\ref{lem:dynamic-async-1} and the spectral bound $\|\tilde r^{(j)}\|_2^2 \leq \lmax \Anorm{\tilde e^{(j)}}^2$. Assembling the coefficients, the telescoped bound takes the form
\begin{equation}
\label{eq:async-telescoped}
E_m \;\leq\; E_0 \;-\; \frac{\nu_{\min}^2\,\tilde\nu_\tau\,\lmin}{n}\,\sum_{j=0}^{m-1} \E[\Anorm{\tilde e^{(j)}}^2],
\end{equation}
where $\tilde\nu_\tau = 1 - 2\rho_{\text{dyn}}\tau$ emerges from the AM--GM-induced combination of the two cross-term contributions in exact analogy with ADG's calculation, with $\rho_{\text{dyn}}$ in place of $\rho$ and an additional $\nu_{\min}^2$ multiplier on the leading progress coefficient. The stated condition $\tau < \sqrt n/(2\nu_{\max}\rho_2)$ guarantees $\tilde\nu_\tau > 0$.

\smallskip\noindent\textit{Step 5: Epoch reduction.} ADG's ``error cannot collapse'' argument now applies, with the worst-case single-step contraction rate set by the \emph{upper} side of Lemma~\ref{lem:dynamic-async-1}: ignoring the non-negative cross terms,
\[
    \E[\Anorm{e^{(j+1)}}^2] \;\geq\; \left(1 - \frac{\nu_{\max}^2\,\lmax}{n}\right) \E[\Anorm{e^{(j)}}^2].
\]
Iterating, the stale errors satisfy $\E[\Anorm{\tilde e^{(j)}}^2] \geq (1/2)\,E_0$ throughout an epoch of length $m \leq \ln(2)\,n/(\nu_{\max}^2\lmax)$. Choosing $m$ at this boundary, the sum is bounded below by the integral of the exponential lower bound on the errors:
\[
    \sum_{j=0}^{m-1}\E[\Anorm{\tilde e^{(j)}}^2] \;\geq\; \frac{n}{2\,\nu_{\max}^2\,\lmax}\,E_0.
\]
Substituting into~\eqref{eq:async-telescoped},
\[
    E_m \;\leq\; E_0 - \frac{\nu_{\min}^2\,\tilde\nu_\tau\,\lmin}{n} \cdot \frac{n}{2\,\nu_{\max}^2\,\lmax}\,E_0
    \;=\; \left(1 - \frac{\nu_{\min}^2}{\nu_{\max}^2}\cdot\frac{\tilde\nu_\tau}{2\kappa}\right) E_0,
\]
where the final equality uses $\lmin/\lmax = 1/\kappa$.
\end{proof}

Both $\nu_{\min}$ and $\nu_{\max}$ appear in the theorem. $\nu_{\min}$
bounds progress from below (imagining the flattest residual) and
multiplies the contraction rate. $\nu_{\max}$ bounds staleness from
above (imagining the spikiest residual) and appears inside
$\rho_{\text{dyn}}$. This pessimism is genuine: when the residual is
spiky the analysis simultaneously (a) underestimates progress using the
flat-residual bound and (b) overestimates staleness using the spiky-residual
bound. See Section~\ref{sec:limitations} for a discussion of this and
other sharpening targets.

The bound in Theorem~\ref{thm:async-power-weighted} differs from the
uniform ADG result in two complementary ways. The leading progress
coefficient $\lmin/n$ is amplified by $\nu_{\min}^2$, an improvement
whenever the residual is non-uniform; the cross-term parameter
$\rho_{\text{dyn}}$ scales as $\nu_{\max}/\sqrt n$ rather than the
uniform $1/n$, a degradation that does not vanish even when the
residual is flat. The two effects pull in opposite directions, and the
ratio $\nu_{\max}/\nu_{\min}$ controls which one dominates. When the
residual is strongly non-uniform along the trajectory (as
Section~\ref{sec:experiments:results} demonstrates is the case for the
problem classes of interest) the progress amplification can
compensate for the looser cross-term, and the bound provides meaningful
information beyond what the uniform ADG result gives.

%\subsection{Comparison with ADG Theorem~2a}

\begin{center}
\renewcommand{\arraystretch}{1.3}
\begin{tabular}{@{}lcc@{}}
\toprule
& \textbf{ADG (uniform)} & \textbf{Dynamic ($\ell=2$)} \\
\midrule
Progress coefficient         & $\lmin/n$                   & $\nu_{\min}^2 \lmin/n$ \\
Cross-term parameter         & $\rho = \|A\|_\infty/n$     & $\rho_{\text{dyn}} = \nu_{\max}\rho_2/\sqrt n$ \\
Scaling (sparse $A$)         & $O(1/n)$                    & $O(\nu_{\max}/\sqrt n)$ \\
Allowable delay              & $\tau = O(n)$               & $\tau = O(\sqrt n/\nu_{\max})$ \\
Epoch length                 & $\geq 0.693\,n/\lmax$       & $\geq 0.693\,n/(\nu_{\max}^2\lmax)$ \\
Epoch reduction              & $1 - \nu_\tau^{\mathrm{ADG}}/(2\kappa)$ & $1 - (\nu_{\min}^2/\nu_{\max}^2)\tilde\nu_\tau/(2\kappa)$ \\
\bottomrule
\end{tabular}
\end{center}

Setting $\nu_{\min} = \nu_{\max} = 1$ (perfectly uniform residual)
recovers the ADG epoch-reduction form $1 - \tilde\nu_\tau/(2\kappa)$
but leaves $\rho_{\text{dyn}} = \rho_2/\sqrt n = O(1/\sqrt n)$. 
%The looseness is entirely due to the Cauchy--Schwarz step.
Theorem~\ref{thm:async-power-weighted} and ADG Theorem~2a should
therefore be compared as complementary results for different sampling
strategies rather than as a generalization and its special case.
When $\nu_{\min}^2 \approx \nu_{\max}^2$ (as observed for all three
problem classes in Section~\ref{sec:experiments} after the initial
transient, with ratios $\nu_{\max}^2/\nu_{\min}^2 \leq 2.4$ on
\texttt{laplace} and $\leq 1.7$ on \texttt{fem}), the per-epoch
reduction $(\nu_{\min}^2/\nu_{\max}^2)\tilde\nu_\tau/(2\kappa)$
approaches the ADG form $\tilde\nu_\tau/(2\kappa)$ over an epoch of
length $\sim n/(\nu_{\max}^2\lmax)$, a factor $\nu_{\max}^2$ shorter
than the ADG epoch. The per-iteration progress is therefore
$\nu_{\max}^2 \cdot (\nu_{\min}^2/\nu_{\max}^2) = \nu_{\min}^2$ times
faster than the uniform baseline. For localized-forcing problems where
$\nu_{\min}^2 \approx 5$, this gives roughly $5\times$ faster
per-iteration convergence than the uniform ADG rate predicts.

Theorem~\ref{thm:async-power-weighted} inherits the consistent-reads
assumption from the ADG framework, which is the standard setting that
makes the stale residual $\tilde r^{(j)}$ a well-defined snapshot
amenable to the conditional-expectation argument used in the proof.
The hardware experiments in Section~\ref{sec:experiments:results} reveal
that this assumption is empirically the \emph{harder} case rather than
the easier one for power-weighted sampling at high concurrency:
consistent reads destabilize the iteration on all three test problems,
while inconsistent reads remain stable, and the underlying mechanism is explored in
Section~\ref{sec:experiments:results}. 
The theorem above should be read as proving a convergence rate
under an idealized read model rather than as recommending consistent
reads for deployment; a rigorous inconsistent-read analog that captures
the feedback-damping effect is the leading direction of follow-up work.

\subsection{Step-Size Control for Large Delays}

The condition $2\rho_{\text{dyn}}\tau < 1$ in Theorem~\ref{thm:async-power-weighted}
can fail if either $\tau$ or $\nu_{\max}$ is large, although the hardware
experiments of Section~\ref{sec:experiments} show the realized
$\rho_{\text{dyn}}\tau$ comfortably inside this window across the full
thread-count sweep. The apparent sharpness is an artifact of fixing the
step size at $\beta = 1$ rather than a feature of the iteration itself:
under-relaxation ($\beta < 1$ in~\eqref{eq:async-update}) recovers
convergence for \emph{any} finite $\tau$ and $\nu_{\max}$, at a rate that
degrades continuously rather than collapsing at a boundary. The condition
demarcates the regime where the full Jacobi step is feasible, not the
regime where convergence is possible at all. This is the dynamic analog
of ADG Theorem~3.

To see this, redo the unrolled recursion of
Theorem~\ref{thm:async-power-weighted} with $\beta$ in place of $1$. The
progress term (from expanding $\Anorm{x^{(j+1)} - x^*}^2$) picks up a
factor $\beta(2-\beta)$, while the cross-coupling damage between stale
and current updates scales as $\beta^2$ since each commit is
$\beta$-scaled. Progress is therefore roughly linear in $\beta$ while
damage is purely quadratic; the asymmetry that makes shrinking
$\beta$ a net win. Repeating the AM--GM decoupling and double-counting
argument with these factors gives the effective convergence parameter
\[
\tilde\nu_\tau(\beta) = 2\beta - \beta^2 - 2\rho_{\text{dyn}}\tau\beta^2
= 2\beta - \beta^2(1 + 2\rho_{\text{dyn}}\tau).
\]
Convergence requires $\tilde\nu_\tau(\beta) > 0$, which rearranges to
$\beta < 2/(1 + 2\rho_{\text{dyn}}\tau)$. This window contains
$\beta = 1$ precisely when $2\rho_{\text{dyn}}\tau < 1$, so
Theorem~\ref{thm:async-power-weighted}'s condition is exactly the
boundary at which the full Jacobi step ceases to be feasible.
Maximizing $\tilde\nu_\tau(\beta)$ by differentiating gives
$\tilde\beta = 1/(1 + 2\rho_{\text{dyn}}\tau)$, which lies in the
window for every $\rho_{\text{dyn}}\tau \geq 0$ and on substitution
yields
\[
\tilde\nu_\tau(\tilde\beta) = \frac{2}{1 + 2\rho_{\text{dyn}}\tau}
- \frac{1 + 2\rho_{\text{dyn}}\tau}{(1 + 2\rho_{\text{dyn}}\tau)^2}
= \frac{1}{1 + 2\rho_{\text{dyn}}\tau}.
\]
The resulting rate degrades smoothly with $\rho_{\text{dyn}}\tau$: at
$\rho_{\text{dyn}}\tau = 0$ we recover the synchronous full-step rate;
as $\rho_{\text{dyn}}\tau \to \infty$ the rate vanishes but never
becomes negative.

\begin{theorem}[Asynchronous power-weighted Jacobi with step-size control]
\label{thm:async-power-weighted-step}
Consider the asynchronous iteration~\eqref{eq:async-update} with step-size $\beta \in (0, 1]$, unit diagonal $A$, power-weighted sampling ($\ell = 2$), consistent reads, and bounded delay $\tau$. Under Definition~\ref{assump:global-ipr}, for any $\tau < \infty$ and any $\nu_{\max} < \infty$, the choice
$$
\tilde\beta = \frac{1}{1 + 2\rho_{\text{dyn}}\tau}
$$
yields convergence with epoch reduction
$$
E_m \leq \left(1 - \frac{\nu_{\min}^2}{2\kappa(1 + 2\rho_{\text{dyn}}\tau)}\right) E_0
$$
for every epoch of length $m \geq 0.693\,n/(\nu_{\min}^2\lmax)$.
\end{theorem}

\begin{proof}
The argument tracks $\beta$-dependence through each step of the proof of Theorem~\ref{thm:async-power-weighted}. The conditional bounds from Step~3 of that proof (Lemma~\ref{lem:dynamic-async-1} and the cross-term inequalities) are statements about expectations over $k_j$ at the stale state and are independent of $\beta$. The $\beta$-dependence enters only through the update~\eqref{eq:async-update} itself. See Appendix~\ref{sect:detailed-proofs} for full details.
\end{proof}

The extension goes through because $\beta$ enters multiplicatively everywhere the component $k_j$ enters. Every term in the unrolled recursion picks up either $\beta$ (from the update itself) or $\beta^2$ (from a cross-coupling), and the AM--GM and double-counting arguments care only about the ratio of these scalings, not about the specific cross-term bound. Any cross-term bound (including the looser dynamic one $\rho_{\text{dyn}}$) slots into the same step-size machinery unchanged.

\subsection{Limitations}
\label{sec:limitations}

While the present analysis establishes a rigorous asynchronous convergence guarantee for residual-weighted sampling, the theoretical bounds remain conservative in several respects. We identify three specific areas where the current framework invites future refinement.

First, the global IPR bound assumes away the direction--delay coupling (Obstacle~1) to provide a worst-case guarantee. However, the IPR trajectory $\nu^2(\tilde r^{(j)})$ is a physical observable whose evolution is demonstrably smooth in practice. Tracking this evolution directly could replace the global assumption with a dynamic, step-wise analysis.

Second, the Cauchy--Schwarz step produces $\nu(\tilde r)/\sqrt n$, losing a factor of $\sqrt n$ relative to ADG. However, because the sampling probabilities $p_r = (\tilde r_r^{(j)})^2/\|\tilde r^{(j)}\|_2^2$ must sum to 1, they inherently form a probability vector with an $\ell_1$ norm of exactly 1. This property naturally invites the use of H\"older's inequality. Bounding the expected matrix entry via $\|p\|_1 \|A_{l,:}\|_\infty$ rather than Cauchy--Schwarz sidesteps the $\ell_2$ penalty entirely, giving the alternative bound $\max_r |A_{lr}| = \|A_{l,:}\|_\infty$. For sparse matrices with bounded entries this is $O(1)$ in $n$ (worse than the ADG $O(1/n)$ but with no $\nu_{\max}$ factor). The right answer is likely a convex combination of Cauchy--Schwarz and H\"older keyed to the residual's R\'enyi-2 entropy, recovering something between $O(1/n)$ and $O(1/\sqrt n)$ depending on concentration.

Third, using $\nu_{\min}$ on progress and $\nu_{\max}$ on damage is strictly worse than tracking $\nu^2(\tilde r^{(j)})$ per step. In practice, spiky residuals cause large progress boosts \emph{at the same moment} that they cause large damage amplifications; the expected product behaves better than the product of expectations. A step-wise analysis could replace the ratio $\nu_{\max}/\nu_{\min}$ with something closer to 1.


\section{Numerical Experiments}
\label{sec:experiments}

The experiments in this section serve two purposes. Synchronous runs on a
small suite of structured problems test the empirical content of the
synchronous theory: whether the IPR
$\nu^2(r) = n\,\|r\|_4^4/\|r\|_2^4$ stays bounded away from its uniform
floor along the trajectory generated by power-weighted sampling, and
whether the resulting sharpened contraction tracks the realized residual
decay. A shared-memory OpenMP implementation then tests the asynchronous
theory under emergent rather than synthetic delays, producing a natural
distribution of delays $\tau$ as a function of thread count from which we
measure both the realized IPR trajectory and the realized delay statistics
directly. We do not inject artificial delays into the shared-memory runs;
the natural delay distribution that arises from concurrent execution is
what we want to measure. We also do not run a separate straggler-resilience
study, viewing that as a distinct question about how the method behaves
outside its analytical regime, which is left as future work.


% \subsection{Test problems}
% \label{sec:experiments:problems}

We use three SPD test problems based on two matrix-free operators, listed
in Table~\ref{tab:problems}. All three are run with the same solver and
instrumentation; the only changes between problem instances are the
operator's apply routine, the diagonal factor used to scale the update,
and the right-hand side. Operators are kept matrix-free to make the
per-update work consistent across thread counts and to avoid having explicit
sparse storage dominate cache behavior at the problem sizes considered
($n$ on the order of $10^4$--$10^5$).

The \texttt{laplace} problem uses the standard five-point Laplacian
stencil on a uniform $N \times N$ grid with Dirichlet boundary conditions
and the smooth right-hand side $b_k = \sin(\pi x_k)\sin(\pi y_k)$. This
is the eigenfunction-like case in which the residual is dominated by
low-frequency components throughout the run, producing long IPR
trajectories with $\nu^2$ modestly above the uniform floor for many
sweeps; it is the hardest of the three problems for the convergence
story and is included to stress the IPR-bound assumption on a residual
that has no localized concentration.

The \texttt{poisson} problem uses the same operator with a localized
point source $b = 100 \cdot e_{N^2/2 + N/2}$ at the grid center. This
is the localized-forcing case where the residual stays spatially
concentrated around the source for many iterations and $\nu^2(r)$
remains elevated well above its uniform floor, giving the sharpened
bound its largest empirical margin over the baseline.

The \texttt{fem} problem uses the consistent mass matrix for
piecewise-linear finite elements on a uniform partition of $[0,1]$
with $n$ nodes and spacing $h = 1/(n-1)$, with a random right-hand
side $b \sim \mathcal{N}(0, I_n)$. The off-diagonals have the opposite
sign convention from the Laplacian ($A_{k,k\pm 1} = h/6 > 0$), the
matrix is well-conditioned, and Jacobi converges rapidly. It serves as
a counterpoint to the Laplacian problems: $\nu^2$ stays close to its
floor, the sharpened bound buys relatively little, and together with
\texttt{poisson} it brackets the range of IPR behavior in structured
problems of interest.

\begin{table}[ht]
\centering
\setlength{\tabcolsep}{3pt}
\begin{tabular}{lccccc}
\hline
Problem        & Operator        & Dim $n$ & $A_{kk}$        & Off-diag         & RHS \\
\hline
\texttt{laplace} & 2D Laplacian   & $N^2$         & $4$             & $-1$             & $\sin(\pi x)\sin(\pi y)$ \\
\texttt{poisson} & 2D Laplacian   & $N^2$         & $4$             & $-1$             & point source (center) \\
\texttt{fem}     & 1D FEM mass    & $N$           & $2h/3$ or $h/3$ & $h/6$            & $\mathcal{N}(0, I_n)$ \\
\hline
\end{tabular}
\caption{Test problems and their matrix-free operator definitions.}
\label{tab:problems}
\end{table}

We report results for $N = 128$ on the 2D problems (so $n = 16{,}384$) and
$n = 8{,}192$ on the FEM problem. These sizes are large enough to be well
inside the sparse-with-bounded-row regime of the ADG analysis and small
enough that several hundred sweeps complete in minutes on a single node.
For problems whose residual is initially concentrated (\texttt{poisson} is
the extreme case, with $\nu^2 = n$ at iteration $0$), the IPR decays
toward a steady-state floor over many sweeps rather than settling quickly,
so unless otherwise noted the IPR statistics reported below are taken
from the second half of $200$-sweep hardware runs, which we have verified
to lie within the asymptotic regime for all three test problems.


% \subsection{Implementation}
% \label{sec:experiments:implementation}

The shared-memory implementation is a single-source C++17 OpenMP
program. Each matrix-free operator exposes a diagonal lookup $A_{kk}$,
an initial-residual routine, and an atomic stencil update applied via
\texttt{\#pragma omp atomic update}; the shared state is a single
iterate $x$ and a single residual $r$, both updated atomically. Each
thread runs an independent loop in which it records a dispatch step
(a relaxed atomic load of a global commit counter), reads the residual,
samples a component $k$ by inverse-CDF sampling on the $\ell^2$-weighted
distribution, computes the update value, and commits via atomic
fetch-and-add on the same counter; the realized delay for that commit
is the commit step minus the dispatch step. Runs were carried out on a
single AMD EPYC node of the ACES cluster at Texas A\&M with
\texttt{OMP\_PROC\_BIND=spread}, \texttt{OMP\_PLACES=cores}, and
exclusive node allocation. At the end of each run the maintained
residual was checked against $b - Ax$ recomputed from the final iterate;
the $\ell^2$ drift was at machine epsilon for every configuration
reported below.
 
A runtime flag selects between two read modes. In
\textbf{consistent-reads} mode, each iteration begins with a
\texttt{memcpy} of the shared residual into a thread-local buffer, and
all subsequent reads of $r$ within the iteration (both the CDF
construction and the lookup of $r_k$ for the update value) come from
that buffer. In \textbf{inconsistent-reads} mode, the per-iteration
\texttt{memcpy} is omitted; every read of $r$ goes through
\texttt{\#pragma omp atomic read} against shared memory, and the read
of $r_k$ used for the update value is performed \emph{a second time}
after sampling. That second read, which may return a different value
than the one used to construct the CDF, is the key inconsistency that
distinguishes this mode from a snapshot-based implementation. The
toggle exists to test empirically whether the consistent-reads
assumption underlying our analysis is a meaningful restriction on the
problems in our test suite, following~\cite{avron2015revisiting}.

We report four sets of measurements: synchronous IPR trajectories and a
comparison of the IPR-sharpened bound (Theorem~\ref{thm:sharp-power})
with realized $A$-norm error; hardware convergence at $T = 96$ threads
under five sampling rules (cyclic, uniform random, and power-weighted
with $\ell \in \{0.5, 1, 2, 4\}$); a thread-count sweep on \texttt{poisson}
from $T=1$ to $T=128$ that probes the IPR's sensitivity to concurrency
and tests Definition~\ref{assump:global-ipr} empirically; and paired
consistent-vs-inconsistent reads at $T = 96$ with matched RNG seeds that
quantify divergence rates. 
%SLURM job scripts and CSV outputs are available in the supplemental material.


\subsection{Results}
\label{sec:experiments:results}

\subsubsection*{Synchronous IPR and bound tightness}

Figure~\ref{fig:combined_ipr} reports synchronous IPR trajectories under
$\ell = 2$ power-weighted sampling on all three test problems, with
numerical values summarized in Table~\ref{tab:ipr-summary}. The IPR sits
well above the uniform-residual floor of $\nu^2 = 1$ throughout every
trajectory. The \texttt{poisson} problem decays from its single-component
initial concentration ($\nu^2 = n$ at iteration $0$) through intermediate
values ($\nu^2 \approx 100, 50, 30, 20$) over the first ten sweeps before
flattening to a steady state around $5.7$, so the bound's progress
amplification is large in the early phase where most of the residual
reduction is achieved and modest in the later phase. Even \texttt{laplace},
which has a smooth right-hand side and no localized forcing, maintains a
sustained IPR near $4.8$: random sampling under power weighting actively
drives the residual toward a non-uniform shape because visited components
are zeroed while unvisited ones retain their initial residual, so even
smooth problems develop sustained concentration once the iteration runs
long enough. The well-conditioned \texttt{fem} case has a similar IPR
range (sustained $\approx 4$) but the trajectory is short because
convergence is rapid.

\begin{table}[ht]
\centering
\small
\begin{tabular}{lccccc}
\toprule
Problem & $n$ & sync $\nu^2_{\min}$ & sync $\nu^2_{\max}$ &
sync ratio $\nu^2_{\max}/\nu^2_{\min}$ & steady-state $\nu^2$ \\
\midrule
\texttt{laplace} & 16{,}384 & 2.22  & 5.43      & 2.4  & $\approx 4.8$ \\
\texttt{poisson} & 16{,}384 & 5.0   & 16{,}384  & --   & $\approx 5.7$ \\
\texttt{fem}     & 8{,}192  & 3.06  & 5.13      & 1.7  & $\approx 4.0$ \\
\bottomrule
\end{tabular}
\caption{Inverse participation ratio statistics across the three test
problems under $\ell = 2$ power-weighted synchronous sampling. The
steady-state value is the mean over the second half of a $200$-sweep
run. For \texttt{poisson}, $\nu^2_{\max}$ is the trivial single-component
bound at iteration $0$; the ratio column is left blank because the
early-iteration $\nu^2$ varies over orders of magnitude and the ratio is
not a meaningful single summary.}
\label{tab:ipr-summary}
\end{table}

\begin{figure}[ht]
    \centering
    \includegraphics[width=\textwidth]{combined_ipr.png}
    \caption{Synchronous IPR trajectories $\nu^2(r^{(i)})$ under
    power-weighted sampling ($\ell = 2$) across the three test problems.
    Steady-state values are roughly $4.8$ for \texttt{laplace}, $5.7$ for
    \texttt{poisson}, and $4.0$ for \texttt{fem}. The \texttt{poisson}
    trajectory decays from its concentrated initial state through
    intermediate values during the first few sweeps before flattening.
    The persistent gap above $1$ is the empirical content of the
    trajectory assumption in Corollary~\ref{cor:global-rate} and
    Theorem~\ref{thm:async-power-weighted}.}
    \label{fig:combined_ipr}
\end{figure}

Figure~\ref{fig:combined_sync_bounds} overlays the simulated error against
the general $\lmin/(n^2 A_{\max})$ bound, the diagonal-weighted bound
$\lmin/(n A_{\max})$, and the IPR-sharpened bound
$\nu^2 \lmin/(n A_{\max})$. The IPR bound is uniformly tighter than the
non-residual-weighted bounds and qualitatively tracks the realized decay
for \texttt{poisson} and \texttt{fem}. For \texttt{laplace} the gap
between the IPR bound and the realized error remains significant,
reflecting that the IPR captures one source of looseness in the analysis
but not all of them; simultaneously tracking the residual's support in
the eigenbasis of $A$ is one natural avenue for further sharpening.

\begin{figure}[ht]
    \centering
    \includegraphics[width=\textwidth]{combined_sync_bounds.png}
    \caption{Theoretical convergence bounds and realized $A$-norm error
    for the three synchronous test problems. The IPR-sharpened bound
    (Theorem~\ref{thm:sharp-power}) sits between the uniform-sampling
    bound (Theorem~\ref{thm:uniform}) and the realized error in every
    case, with the gap most pronounced for \texttt{poisson} where the
    sustained IPR is largest. The simulated error decays substantially
    faster than the sharpened bound, indicating that further sharpening
    is possible.}
    \label{fig:combined_sync_bounds}
\end{figure}

\subsubsection*{Hardware convergence and IPR stability under concurrency}

Under emergent rather than synthetic delays, power-weighted sampling
translates the IPR advantage into a measurable convergence improvement.
Figure~\ref{fig:hw_convergence} reports residual decay under five sampling
rules at $T = 96$ threads on a single ACES node. The power-weighted rules
($\ell \in \{1, 2, 4\}$) reach residual $10^{-3}$ on \texttt{poisson} in
roughly half the sweeps required by uniform random or cyclic sampling;
the \texttt{laplace} advantage is smaller but consistent, matching the
smaller IPR of that problem; and \texttt{fem} converges so rapidly under
all rules that the comparison is dominated by start-up effects. The
near-coincidence of $\ell = 2$ and $\ell = 4$ curves suggests diminishing
returns from increasing residual bias once concentration is fully exploited.

\begin{figure}[ht]
    \centering
    \includegraphics[width=\textwidth]{fig_hw_convergence.png}
    \caption{Hardware convergence under five sampling rules at $T=96$ threads,
    inconsistent reads. Power-weighted sampling with $\ell \in \{1, 2, 4\}$
    converges measurably faster than uniform random or cyclic sampling on
    a per-update basis, with the largest advantage for \texttt{poisson}.}
    \label{fig:hw_convergence}
\end{figure}

The hardware IPR trajectories in Figure~\ref{fig:hw_ipr} closely follow
their synchronous counterparts, providing direct empirical evidence that
the global trajectory bound in Definition~\ref{assump:global-ipr} is not
violated by hardware-induced delays. Uniform and cyclic sampling produce
noticeably \emph{higher} IPR values on \texttt{poisson} and \texttt{fem}
because they fail to depress the high-residual components fast enough:
non-power-weighted rules leave concentration on the table.

\begin{figure}[ht]
    \centering
    \includegraphics[width=\textwidth]{fig_hw_ipr.png}
    \caption{Hardware IPR trajectories under the configuration of
    Figure~\ref{fig:hw_convergence}. Power-weighted sampling tracks the
    synchronous trajectories of Figure~\ref{fig:combined_ipr}, confirming
    that residual concentration is preserved under emergent asynchronous
    execution. Uniform and cyclic sampling produce substantially higher
    IPR for \texttt{poisson} and \texttt{fem}, reflecting that those
    rules visit components less aggressively in proportion to residual
    mass.}
    \label{fig:hw_ipr}
\end{figure}

The thread-count sweep on \texttt{poisson} (Figure~\ref{fig:hw_scaling}
and Table~\ref{tab:hw-scaling}) gives three findings. Wall time scales
near-linearly up to $\sim 64$ threads and flattens beyond, reflecting
memory bandwidth and atomic contention on the EPYC node rather than any
algorithmic limit. Realized mean and maximum delay both grow sub-linearly
with thread count, and at $T=128$ the realized mean delay is well inside
the convergence window predicted by
Theorem~\ref{thm:async-power-weighted}. Most importantly for the theory,
the sustained $\nu^2$ varies by less than $0.2$ in absolute terms
($\sim 3\%$) across the entire sweep from $T = 1$ to $T = 128$. This is
the strongest empirical support for Definition~\ref{assump:global-ipr}:
even under heavy concurrent execution with realized mean delays in the
$10^2$ commit range, the residual concentration remains essentially
identical to the synchronous baseline, justifying treating the IPR as a
structural property of the iteration trajectory rather than an artifact
of the delay model.

\begin{figure}[ht]
    \centering
    \includegraphics[width=\textwidth]{fig_hw_scaling.png}
    \caption{Thread-count sweep for \texttt{poisson}, power-weighted
    ($\ell=2$), inconsistent reads. (a) Wall time scales near-linearly up
    to $\sim 64$ threads before flattening. (b) Realized mean and maximum
    delay both grow sub-linearly with thread count; the $\tau=1$ reference
    line places the realized regime well inside the convergence window of
    Theorem~\ref{thm:async-power-weighted}. (c) The sustained IPR is
    essentially independent of thread count, validating
    Definition~\ref{assump:global-ipr} empirically. An exploration of the occasional divergence for $T=96$ is given in
    Table~\ref{tab:consistency-sweep}.}
    \label{fig:hw_scaling}
\end{figure}

\begin{table}[ht]
\centering
\small
\renewcommand{\arraystretch}{1.1}
\begin{tabular}{r r r r r}
\toprule
$T$ & Wall time (s) & $\bar\tau$ & $\tau_{\max}$ & Sustained $\nu^2$ \\
\midrule
  1  &   70.05 &     0.0 &         0 & 5.67 \\
  2  &   37.57 &     1.0 &         6 & 5.67 \\
  4  &   18.66 &     3.0 &        65 & 5.67 \\
  8  &    9.46 &     6.9 &        42 & 5.67 \\
 16  &    4.89 &    14.8 &       208 & 5.68 \\
 32  &    2.60 &    29.8 &    1\,368 & 5.70 \\
 64  &    1.38 &    59.5 &    9\,667 & 5.73 \\
 96\tablefootnote{Averaged over the converged subset of 6 jobs at this thread count; see Table 4 for the corresponding divergence-rate study with seed variation.}  &    0.95 &    89.5 &    6\,928 & 5.74 \\
128  &    1.02 &   111.6 &   89\,607 & 5.84 \\
\bottomrule
\end{tabular}
\caption{Hardware scaling for \texttt{poisson}, $N=128$, power-weighted
($\ell=2$), inconsistent reads. $\bar\tau$ and
$\tau_{\max}$ are the mean and cumulative-maximum commit-delay over the
second half of the run; sustained $\nu^2$ is the mean IPR over the same
window.}
\label{tab:hw-scaling}
\end{table}

\subsubsection*{Consistent vs.\ inconsistent reads}

The original ADG analysis treats consistent reads as the easy case for
the convergence theory and inconsistent reads as a more delicate setting
that requires a separate analysis (\cite[Section~7]{avron2015revisiting}).
Under power-weighted sampling on a concentrated residual, the empirical
ordering is reversed. Table~\ref{tab:consistency-sweep} reports a paired
sweep on \texttt{poisson} at $T = 96$ with six matched seeds: inconsistent
reads diverged in 1 of 6 runs (the same residual-collision phenomenon
documented in Table~\ref{tab:hw-scaling}), while consistent reads diverged
in all 6. Combined with six earlier consistent runs at seed=42 that also
all diverged, the divergence rate for consistent reads at this
configuration is $12/12$ versus $1/6$ for inconsistent reads with proper
seed variation. The per-iteration cost of consistent reads exceeded that
of inconsistent reads on every problem by $10$--$25\%$ due to the
iteration-start \texttt{memcpy}; on \texttt{fem}, where inconsistent
reads reach the $10^{-6}$ tolerance within roughly $30$ sweeps, the
consistent runs ran to the $200$-sweep cap with a final wall time
$25\times$ higher.

The mechanism we conjecture is that the IPR plays a dual role under
power weighting: it amplifies expected progress
(Theorem~\ref{thm:sharp-power}) but also amplifies the collision rate
between concurrent threads. The argument in~\cite{avron2015revisiting}
that inconsistent-read violations are rare in the sparse reference
scenario relies on each update landing on a uniformly random component:
another thread's update hits a fixed read set $S$ with probability at
most $|S|/n$. Under power weighting with $\ell = 2$ the per-component
update probabilities are $r_k^2 / \|r\|_2^2$ rather than $1/n$, and a
short calculation using $\|r\|_\infty^4 \leq \|r\|_4^4 = \nu^2 \|r\|_2^4 / n$
shows the collision-probability bound becomes
$|S| \cdot \sqrt{\nu^2/n}$. The per-iteration collision probability is
therefore inflated by a factor at most $\sqrt{n\nu^2}$ relative to the
uniform case; roughly $305\times$ at the parameters of our \texttt{poisson}
runs. The expected number of concurrent thread collisions on a single
component per sampling round is $\binom{T}{2}\nu^2/n$; even modest IPR
($\nu^2 \approx 3$ on \texttt{laplace}) inflates the collision rate by
orders of magnitude relative to uniform sampling. 

Using consistent
reads, every such collision contributes a frozen-value pile-on whose net
effect is to overshoot the residual at the colliding component by a
factor proportional to the number of colliding threads, with the
\texttt{poisson} initial $\nu^2 \approx n$ producing catastrophic
overshoot in a single sweep and the smaller \texttt{laplace} and
\texttt{fem} rates producing gradual accumulation that nonetheless
reaches divergence within a few thousand sweeps. Under inconsistent
reads, the second residual read used to set the update value sees the
partial effect of any in-flight commits to the colliding component and
shrinks the local update accordingly: a sign-aware correction that the
absolute-value bound above cannot capture. The same $\nu^2$ that
amplifies progress in Theorem~\ref{thm:sharp-power} also amplifies
collision probability ($\nu^2/n$ per pair); under consistent reads
these roles compound destructively, while under inconsistent reads the
feedback mechanism mediates between them. A rigorous inconsistent-read
analog of Theorem~\ref{thm:async-power-weighted} that captures this
distinction (rather than just bounding collision frequency from above)
is the natural next theoretical step, requiring an analysis that tracks
the sign-correlation between the stale read and the in-flight update
which the current proof architecture handles only as a worst-case
absolute value.

\begin{table}[ht]
\centering
\small
\renewcommand{\arraystretch}{1.1}
\begin{tabular}{l r r c c}
\toprule
Problem & $n$ & $T$ & Inconsistent & Consistent \\
\midrule
\multicolumn{5}{l}{\textit{Original sweep (baseline $n$)}}\\
\texttt{poisson} & $16{,}384$  & 96 & 5/6 conv. & 0/6 conv. \\
\texttt{laplace} & $16{,}384$  & 96 & 6/6 stable & 0/6 conv. \\
\texttt{fem}     & $8{,}192$   & 96 & 6/6 conv.  & 0/6 conv. \\
\midrule
\multicolumn{5}{l}{\textit{Scaling test ($4\times$ larger $n$)}}\\
\texttt{poisson} & $65{,}536$  & 96 & 6/6 conv.  & 0/6 conv. \\
\texttt{laplace} & $65{,}536$  & 96 & 6/6 stable & \textbf{6/6 conv.} \\
\texttt{fem}     & $32{,}768$  & 96 & 6/6 conv.  & \textbf{6/6 conv.} \\
\bottomrule
\end{tabular}
\caption{Paired consistent-vs-inconsistent runs across the three test
problems, power-weighted ($\ell=2$), $T=96$, with six matched seeds.
At baseline $n$ all three problems show $0/6$ consistent-reads convergence.
Quadrupling $n$ rescues consistent reads on \texttt{laplace} and
\texttt{fem} ($6/6$ in both cases), because the per-pair collision
probability $\nu^2/n$ scales as $1/n$ when $\nu^2$ is bounded. On
\texttt{poisson} the localized point-source forcing produces
$\nu^2 \propto n$ initially, so the collision probability is
$n$-independent and quadrupling $n$ does not help.}
\label{tab:consistency-sweep}
\end{table}

\begin{table}[ht]
\centering
\small
\renewcommand{\arraystretch}{1.1}
\begin{tabular}{r c c c c c}
\toprule
$T$ & 8 & 16 & 32 & 64 & 96 \\
\midrule
Inconsistent (conv.) & 6/6 & 6/6 & 6/6 & 6/6 & 5/6 \\
Consistent  (conv.)  & 0/6 & 0/6 & 0/6 & 0/6 & 0/6 \\
\bottomrule
\end{tabular}
\caption{Concurrency sweep on \texttt{poisson} ($N=128$, $\ell=2$).
Consistent reads diverges at every tested thread count: the iteration-0
pile-on at the point-source residual is catastrophic for any $T \geq 2$
because the entire sampling distribution collapses onto a single component
while all $T$ threads see the same snapshot value. Inconsistent reads
remains stable across the entire range.}
\label{tab:threshold-sweep}
\end{table}

The four experimental panels point to a single mechanism: the IPR
controls both how much power-weighted sampling helps in expectation and
how much it hurts under concurrency. On problems where $\nu^2$ stays
bounded as $n$ grows (\texttt{laplace}, \texttt{fem}), both the progress
amplification and the collision rate scale modestly, and quadrupling $n$
is enough to rescue consistent reads (Table~\ref{tab:consistency-sweep}).
On problems where $\nu^2$ scales with $n$ (\texttt{poisson} with its
point-source forcing), the progress amplification grows linearly with
problem size but the collision rate stays roughly constant, so consistent
reads remains broken at every $n$ we tested. In both regimes,
inconsistent reads provides the sign-aware feedback needed to absorb the
collisions. The practical recommendation is to default to inconsistent
reads whenever power-weighted sampling is used at non-trivial thread
count, and to treat the IPR not as a side observable but as the
structural quantity that simultaneously determines the algorithm's
productivity and its safety.

%===========================================================================
\section{Discussion and Conclusion}
\label{sec:conclusion}
%===========================================================================

The synchronous analysis presented here shows that the standard convergence
bound for residual-weighted randomized Jacobi is systematically loose by a
factor equal to the inverse participation ratio of the residual.
Theorem~\ref{thm:sharp-power} makes this precise: the per-step progress is
amplified by exactly $\nu^2$, a scalar that is always at least $1$, that is
strictly greater than $1$ whenever the residual is non-uniform, and that
can be computed from the residual at essentially no additional cost. The
numerical experiments in Section~\ref{sec:experiments} show that this
sharpened bound is not a theoretical curiosity. The IPR is sustained in
the range $4$--$6$ across all three test problem classes with substantially
higher values during the early-iteration phase where most of the residual
reduction happens, and it grows along the trajectory before stabilizing
because random sampling drives visited components toward their local
solutions while unvisited components retain their initial residual.

The asynchronous extension carries this analysis into the
Avron--Druinsky--Gupta framework under a global bound on the IPR
trajectory, but not cleanly: the Cauchy--Schwarz step in the cross-term
bound introduces a $\sqrt{n}$ scaling loss, so
Theorem~\ref{thm:async-power-weighted} does not subsume ADG Theorem~2a as
a special case. The two results are best read as complementary bounds
for different sampling strategies. Hardware experiments on the ACES
cluster show that the IPR advantage translates into measurably faster
convergence under real concurrent execution
(Figure~\ref{fig:hw_convergence}), that the IPR trajectory is nearly
insensitive to thread count across two decades of concurrency
(Figure~\ref{fig:hw_scaling}, varying by less than $3\%$ on
\texttt{poisson} from $T=1$ to $T=128$ despite mean delays growing by two
orders of magnitude), and that the realized delays sit well inside the
convergence window. The $\sqrt{n}$ cross-term looseness therefore does
not appear binding in practice for the problem classes considered, even
though the theoretical gap remains.

The least anticipated finding is the reversal of the standard ADG
narrative: consistent reads, the easy case for the convergence theory,
is empirically the unsafe case for power-weighted sampling at high
concurrency, while inconsistent reads remains stable. The mechanism
proposed in Section~\ref{sec:experiments:results} ties this directly to
the IPR: the same $\nu^2$ that amplifies expected progress in
Theorem~\ref{thm:sharp-power} also amplifies the per-pair thread
collision rate by a factor of $\nu^2$ over the uniform case. Under
consistent reads, collisions compound into pile-up overshoots; under
inconsistent reads, the second residual read provides a sign-aware
correction. The two roles of $\nu^2$ make a sharp empirical prediction
that we are now able to verify directly. On problems where $\nu^2$ stays
bounded as $n$ grows, both the speedup and the per-pair collision rate
$\nu^2/n$ scale modestly, so consistent reads should become safe in the
asymptotic regime; Table~\ref{tab:consistency-sweep} confirms this,
with quadrupling $n$ taking \texttt{laplace} and \texttt{fem} from
$0/6$ converged to $6/6$ under consistent reads. On problems where
$\nu^2$ scales with $n$ (e.g., localized forcing whose residual
concentration is a feature of the physics), the speedup grows linearly
with $n$ while the collision rate stays constant, and consistent reads
fails at every scale; \texttt{poisson} at $n = 65{,}536$ still diverges
in all six trials.

Several directions follow naturally. The most pressing is a rigorous
inconsistent-read analog of Theorem~\ref{thm:async-power-weighted} that
captures the feedback-damping effect, which would require tracking the
sign-correlation between the stale read and the in-flight update rather
than treating the cross-term as a worst-case absolute value. Given that
inconsistent reads turn out to be strictly preferable to consistent
reads for the problem class of interest, this is the direction that most
directly informs practitioner choice. A natural-seeming algorithmic
mitigation, preceding power-weighted sampling with $K$ uniform warmup
sweeps so that the residual diffuses before residual-weighted sampling
activates, turns out empirically not to help: a scan of
$K \in \{0, 5, 10, 20, 50, 100, 150\}$ on \texttt{poisson} at $T = 96$
gave consistent-reads divergence in $42/42$ trials regardless of $K$.
The pile-on is self-reinforcing rather than transient, because $T$
concurrent commits to any high-residual component create a fresh hot
spot that attracts the next round of power-weighted samples.

A second direction is to replace the global IPR bound in
Definition~\ref{assump:global-ipr} with a martingale or filtration
argument that tracks $\nu^2(\tilde r^{(j)})$ directly. The empirical
observations that the IPR evolves smoothly rather than fluctuating
wildly and that it is concurrency-insensitive provide scaffolding for
such an analysis; the analytical task is to convert that smoothness
into a step-to-step recurrence. A related direction is to prove
$\nu_i^2 \geq \bar\nu^2 > 1$ analytically for specific matrix classes
such as discretized elliptic PDEs with localized forcing, which would
convert Corollary~\ref{cor:global-rate} from a conditional statement
into an unconditional one for the problem class where the sharpened
bound shows the largest empirical margin. Finally, the analysis here
treats only $\ell = 2$ and could be extended to a general-$\ell$ family
of non-uniformity parameters, with the optimal $\ell$ as a function of
the residual distribution as the natural object of study and a clear
connection to the sketch-and-project framework via what residual-weighted
sketch selection does to $\lambda_{\min}(\E[Z])$.

\bmhead{Acknowledgements}

This work leveraged the ACES Cluster at Texas A\&M University under
allocation CIS250436 from the Advanced Cyberinfrastructure Coordination
Ecosystem: Services \& Support (ACCESS) program, which is supported by
U.S. National Science Foundation grants \#2138259, \#2138286, \#2138307,
\#2137603, and \#2138296.

% \section*{Declarations}

% Some journals require declarations to be submitted in a standardised format. Please check the Instructions for Authors of the journal to which you are submitting to see if you need to complete this section. If yes, your manuscript must contain the following sections under the heading `Declarations':

% \begin{itemize}
% \item Funding
% \item Conflict of interest/Competing interests (check journal-specific guidelines for which heading to use)
% \item Ethics approval and consent to participate
% \item Consent for publication
% \item Data availability 
% \item Materials availability
% \item Code availability 
% \item Author contribution
% \end{itemize}

% \noindent
% If any of the sections are not relevant to your manuscript, please include the heading and write `Not applicable' for that section. 



\begin{appendices}

\section{Detailed Proof of Theorem~\ref{thm:async-power-weighted-step}}
\label{sect:detailed-proofs}
\begin{proof}
\label{proof:detailed-step-size}
\noindent\textit{Step 1: $\beta$-modified one-step identity.} With
$\beta \in (0, 1]$, the update gives
$e^{(j+1)} = e^{(j)} + \beta\,\tilde r_{k_j}^{(j)} e_{k_j}$. Each
committed update in the in-flight residual decomposition now carries an
additional factor of $\beta$:
\[
r^{(j)} = \tilde r^{(j)} - \beta \sum_{l \in J_j} \tilde r_{k_l}^{(l)} A e_{k_l}.
\]
Expanding $\Anorm{e^{(j+1)}}^2$ with these factors,
\begin{equation}
\label{eq:async-onestep-beta}
\Anorm{e^{(j+1)}}^2 = \Anorm{e^{(j)}}^2 - \beta(2-\beta)\,(\tilde r_{k_j}^{(j)})^2 + 2\beta^2 \sum_{l \in J_j} A_{k_j, k_l}\,\tilde r_{k_j}^{(j)}\,\tilde r_{k_l}^{(l)}.
\end{equation}
Compared to the $\beta = 1$ identity~\eqref{eq:async-onestep-proof}, the
progress coefficient scales as $\beta(2-\beta)$ (from the $-2\beta + \beta^2$
combination in the expansion) and the interference scales as $\beta^2$
(one factor of $\beta$ from the current update and one from the in-flight
update being measured against it).

\smallskip\noindent\textit{Step 2: Conditional expectations and
telescoping.} Steps~2--4 of the proof of
Theorem~\ref{thm:async-power-weighted} apply with the same conditional
bounds, modulated only by the $\beta$ scalings
from~\eqref{eq:async-onestep-beta}. The telescoped bound becomes
\[
E_m \;\leq\; E_0 \;-\; \frac{\nu_{\min}^2\,\lmin}{n}\,\tilde\nu_\tau(\beta)\,\sum_{j=0}^{m-1} \E[\Anorm{\tilde e^{(j)}}^2],
\]
where the effective convergence parameter combines the progress and
interference $\beta$-scalings:
\[
\tilde\nu_\tau(\beta) \;:=\; \beta(2-\beta) - 2\beta^2\rho_{\text{dyn}}\tau \;=\; 2\beta - \beta^2(1 + 2\rho_{\text{dyn}}\tau).
\]
Convergence requires $\tilde\nu_\tau(\beta) > 0$, which is equivalent to
$\beta < 2/(1 + 2\rho_{\text{dyn}}\tau)$. Unlike the $\beta = 1$ case,
no upper bound on $\rho_{\text{dyn}}\tau$ is required: for any finite
delay and any finite IPR bound, the convergence window for $\beta$ is
non-empty.

\smallskip\noindent\textit{Step 3: Optimal step size.} Maximizing
$\tilde\nu_\tau(\beta)$ over $\beta \in (0, 2/(1 + 2\rho_{\text{dyn}}\tau))$:
\[
\frac{d}{d\beta}\tilde\nu_\tau(\beta) = 2 - 2\beta\,(1 + 2\rho_{\text{dyn}}\tau) = 0 \quad\Longrightarrow\quad \tilde\beta = \frac{1}{1 + 2\rho_{\text{dyn}}\tau}.
\]
The optimizer $\tilde\beta$ lies in $(0, 1]$ for every
$\rho_{\text{dyn}}\tau \geq 0$, so the choice is always feasible (and
reduces to $\beta = 1$ in the synchronous limit $\tau = 0$). Substituting:
\[
\tilde\nu_\tau(\tilde\beta) = \frac{2}{1 + 2\rho_{\text{dyn}}\tau} - \frac{1 + 2\rho_{\text{dyn}}\tau}{(1 + 2\rho_{\text{dyn}}\tau)^2} = \frac{1}{1 + 2\rho_{\text{dyn}}\tau}.
\]

\smallskip\noindent\textit{Step 4: Epoch reduction.} The epoch-assembly
argument from Step~5 of the proof of
Theorem~\ref{thm:async-power-weighted} is unchanged (it depends only on
Lemma~\ref{lem:dynamic-async-1}, which is $\beta$-free). Substituting
the optimal $\tilde\nu_\tau(\tilde\beta)$ into the epoch bound:
\[
E_m \;\leq\; \left(1 - \frac{\nu_{\min}^2\,\tilde\nu_\tau(\tilde\beta)}{2\kappa}\right) E_0 \;=\; \left(1 - \frac{\nu_{\min}^2}{2\kappa(1 + 2\rho_{\text{dyn}}\tau)}\right) E_0.
\]
\end{proof}

\end{appendices}

%%===========================================================================================%%
%% If you are submitting to one of the Nature Portfolio journals, using the eJP submission   %%
%% system, please include the references within the manuscript file itself. You may do this  %%
%% by copying the reference list from your .bbl file, paste it into the main manuscript .tex %%
%% file, and delete the associated \verb+\bibliography+ commands.                            %%
%%===========================================================================================%%

\bibliography{references}

\end{document}
